\documentclass[11pt,a4paper]{article}

\usepackage[T1]{fontenc}
\usepackage[utf8]{inputenc}
\usepackage{lmodern}
\usepackage[margin=2.4cm]{geometry}
\usepackage{amsmath,amssymb}
\usepackage{booktabs}
\usepackage{longtable}
\usepackage{array}
\usepackage{graphicx}
\usepackage{pgfplots}
\pgfplotsset{compat=1.18}
\usepackage{microtype}
\usepackage[hidelinks]{hyperref}
\usepackage{url}
\usepackage[numbers,sort&compress]{natbib}
\setlength{\emergencystretch}{2em}

% tex/numbers.tex is written by `papertool ledger` from data/measurements.json. Every
% quantitative claim below goes through one of these lookups; `papertool numbers` fails the
% build on an unknown key, so a number that no experiment produced cannot be printed. An
% undefined key expands to a visible marker rather than to nothing.
% Generated by `papertool ledger` from data/measurements.json -- do not edit.
% Every quantitative claim in the paper resolves through one of these names.
\makeatletter
\expandafter\def\csname pv@margolus_size8.x\endcsname{8}
\expandafter\def\csname pv@margolus_size8.peak_ws_mb\endcsname{33.16}
\expandafter\def\csname pv@margolus_size8.peak_ws_sd_mb\endcsname{0.2132}
\expandafter\def\csname pv@margolus_size8.elapsed_s\endcsname{0.2899}
\expandafter\def\csname pv@margolus_size8.bit_exact_reversal\endcsname{10/10}
\expandafter\def\csname pvfrac@margolus_size8.bit_exact_reversal\endcsname{1}
\expandafter\def\csname pvn@margolus_size8.bit_exact_reversal\endcsname{10}
\expandafter\def\csname pvtrue@margolus_size8.bit_exact_reversal\endcsname{10}
\expandafter\def\csname pv@margolus_size8.endpoint_recovered\endcsname{10/10}
\expandafter\def\csname pvfrac@margolus_size8.endpoint_recovered\endcsname{1}
\expandafter\def\csname pvn@margolus_size8.endpoint_recovered\endcsname{10}
\expandafter\def\csname pvtrue@margolus_size8.endpoint_recovered\endcsname{10}
\expandafter\def\csname pv@margolus_size8.entropy_conserved\endcsname{10/10}
\expandafter\def\csname pvfrac@margolus_size8.entropy_conserved\endcsname{1}
\expandafter\def\csname pvn@margolus_size8.entropy_conserved\endcsname{10}
\expandafter\def\csname pvtrue@margolus_size8.entropy_conserved\endcsname{10}
\expandafter\def\csname pv@margolus_size8.hamming_weight\endcsname{11.9}
\expandafter\def\csname pvci@margolus_size8.hamming_weight\endcsname{11.9 (2.132)}
\expandafter\def\csname pvsd@margolus_size8.hamming_weight\endcsname{2.132}
\expandafter\def\csname pvn@margolus_size8.hamming_weight\endcsname{10}
\expandafter\def\csname pvmin@margolus_size8.hamming_weight\endcsname{7}
\expandafter\def\csname pvmax@margolus_size8.hamming_weight\endcsname{14}
\expandafter\def\csname pv@margolus_size8.midpoint_match\endcsname{10/10}
\expandafter\def\csname pvfrac@margolus_size8.midpoint_match\endcsname{1}
\expandafter\def\csname pvn@margolus_size8.midpoint_match\endcsname{10}
\expandafter\def\csname pvtrue@margolus_size8.midpoint_match\endcsname{10}
\expandafter\def\csname pv@margolus_size16.x\endcsname{16}
\expandafter\def\csname pv@margolus_size16.peak_ws_mb\endcsname{33.18}
\expandafter\def\csname pv@margolus_size16.peak_ws_sd_mb\endcsname{0.1905}
\expandafter\def\csname pv@margolus_size16.elapsed_s\endcsname{0.3003}
\expandafter\def\csname pv@margolus_size16.bit_exact_reversal\endcsname{10/10}
\expandafter\def\csname pvfrac@margolus_size16.bit_exact_reversal\endcsname{1}
\expandafter\def\csname pvn@margolus_size16.bit_exact_reversal\endcsname{10}
\expandafter\def\csname pvtrue@margolus_size16.bit_exact_reversal\endcsname{10}
\expandafter\def\csname pv@margolus_size16.endpoint_recovered\endcsname{10/10}
\expandafter\def\csname pvfrac@margolus_size16.endpoint_recovered\endcsname{1}
\expandafter\def\csname pvn@margolus_size16.endpoint_recovered\endcsname{10}
\expandafter\def\csname pvtrue@margolus_size16.endpoint_recovered\endcsname{10}
\expandafter\def\csname pv@margolus_size16.entropy_conserved\endcsname{10/10}
\expandafter\def\csname pvfrac@margolus_size16.entropy_conserved\endcsname{1}
\expandafter\def\csname pvn@margolus_size16.entropy_conserved\endcsname{10}
\expandafter\def\csname pvtrue@margolus_size16.entropy_conserved\endcsname{10}
\expandafter\def\csname pv@margolus_size16.hamming_weight\endcsname{47.7}
\expandafter\def\csname pvci@margolus_size16.hamming_weight\endcsname{47.7 (5.559)}
\expandafter\def\csname pvsd@margolus_size16.hamming_weight\endcsname{5.559}
\expandafter\def\csname pvn@margolus_size16.hamming_weight\endcsname{10}
\expandafter\def\csname pvmin@margolus_size16.hamming_weight\endcsname{37}
\expandafter\def\csname pvmax@margolus_size16.hamming_weight\endcsname{55}
\expandafter\def\csname pv@margolus_size16.midpoint_match\endcsname{10/10}
\expandafter\def\csname pvfrac@margolus_size16.midpoint_match\endcsname{1}
\expandafter\def\csname pvn@margolus_size16.midpoint_match\endcsname{10}
\expandafter\def\csname pvtrue@margolus_size16.midpoint_match\endcsname{10}
\expandafter\def\csname pv@margolus_size32.x\endcsname{32}
\expandafter\def\csname pv@margolus_size32.peak_ws_mb\endcsname{33.06}
\expandafter\def\csname pv@margolus_size32.peak_ws_sd_mb\endcsname{0.1444}
\expandafter\def\csname pv@margolus_size32.elapsed_s\endcsname{0.2552}
\expandafter\def\csname pv@margolus_size32.bit_exact_reversal\endcsname{10/10}
\expandafter\def\csname pvfrac@margolus_size32.bit_exact_reversal\endcsname{1}
\expandafter\def\csname pvn@margolus_size32.bit_exact_reversal\endcsname{10}
\expandafter\def\csname pvtrue@margolus_size32.bit_exact_reversal\endcsname{10}
\expandafter\def\csname pv@margolus_size32.endpoint_recovered\endcsname{10/10}
\expandafter\def\csname pvfrac@margolus_size32.endpoint_recovered\endcsname{1}
\expandafter\def\csname pvn@margolus_size32.endpoint_recovered\endcsname{10}
\expandafter\def\csname pvtrue@margolus_size32.endpoint_recovered\endcsname{10}
\expandafter\def\csname pv@margolus_size32.entropy_conserved\endcsname{10/10}
\expandafter\def\csname pvfrac@margolus_size32.entropy_conserved\endcsname{1}
\expandafter\def\csname pvn@margolus_size32.entropy_conserved\endcsname{10}
\expandafter\def\csname pvtrue@margolus_size32.entropy_conserved\endcsname{10}
\expandafter\def\csname pv@margolus_size32.hamming_weight\endcsname{203.3}
\expandafter\def\csname pvci@margolus_size32.hamming_weight\endcsname{203.3 (9.889)}
\expandafter\def\csname pvsd@margolus_size32.hamming_weight\endcsname{9.889}
\expandafter\def\csname pvn@margolus_size32.hamming_weight\endcsname{10}
\expandafter\def\csname pvmin@margolus_size32.hamming_weight\endcsname{190}
\expandafter\def\csname pvmax@margolus_size32.hamming_weight\endcsname{219}
\expandafter\def\csname pv@margolus_size32.midpoint_match\endcsname{10/10}
\expandafter\def\csname pvfrac@margolus_size32.midpoint_match\endcsname{1}
\expandafter\def\csname pvn@margolus_size32.midpoint_match\endcsname{10}
\expandafter\def\csname pvtrue@margolus_size32.midpoint_match\endcsname{10}
\expandafter\def\csname pv@margolus_size48.x\endcsname{48}
\expandafter\def\csname pv@margolus_size48.peak_ws_mb\endcsname{33.27}
\expandafter\def\csname pv@margolus_size48.peak_ws_sd_mb\endcsname{0.1959}
\expandafter\def\csname pv@margolus_size48.elapsed_s\endcsname{0.2508}
\expandafter\def\csname pv@margolus_size48.bit_exact_reversal\endcsname{10/10}
\expandafter\def\csname pvfrac@margolus_size48.bit_exact_reversal\endcsname{1}
\expandafter\def\csname pvn@margolus_size48.bit_exact_reversal\endcsname{10}
\expandafter\def\csname pvtrue@margolus_size48.bit_exact_reversal\endcsname{10}
\expandafter\def\csname pv@margolus_size48.endpoint_recovered\endcsname{10/10}
\expandafter\def\csname pvfrac@margolus_size48.endpoint_recovered\endcsname{1}
\expandafter\def\csname pvn@margolus_size48.endpoint_recovered\endcsname{10}
\expandafter\def\csname pvtrue@margolus_size48.endpoint_recovered\endcsname{10}
\expandafter\def\csname pv@margolus_size48.entropy_conserved\endcsname{10/10}
\expandafter\def\csname pvfrac@margolus_size48.entropy_conserved\endcsname{1}
\expandafter\def\csname pvn@margolus_size48.entropy_conserved\endcsname{10}
\expandafter\def\csname pvtrue@margolus_size48.entropy_conserved\endcsname{10}
\expandafter\def\csname pv@margolus_size48.hamming_weight\endcsname{459.3}
\expandafter\def\csname pvci@margolus_size48.hamming_weight\endcsname{459.3 (24.01)}
\expandafter\def\csname pvsd@margolus_size48.hamming_weight\endcsname{24.01}
\expandafter\def\csname pvn@margolus_size48.hamming_weight\endcsname{10}
\expandafter\def\csname pvmin@margolus_size48.hamming_weight\endcsname{424}
\expandafter\def\csname pvmax@margolus_size48.hamming_weight\endcsname{495}
\expandafter\def\csname pv@margolus_size48.midpoint_match\endcsname{10/10}
\expandafter\def\csname pvfrac@margolus_size48.midpoint_match\endcsname{1}
\expandafter\def\csname pvn@margolus_size48.midpoint_match\endcsname{10}
\expandafter\def\csname pvtrue@margolus_size48.midpoint_match\endcsname{10}
\expandafter\def\csname pv@margolus_size64.x\endcsname{64}
\expandafter\def\csname pv@margolus_size64.peak_ws_mb\endcsname{33.28}
\expandafter\def\csname pv@margolus_size64.peak_ws_sd_mb\endcsname{0.1676}
\expandafter\def\csname pv@margolus_size64.elapsed_s\endcsname{0.2574}
\expandafter\def\csname pv@margolus_size64.bit_exact_reversal\endcsname{10/10}
\expandafter\def\csname pvfrac@margolus_size64.bit_exact_reversal\endcsname{1}
\expandafter\def\csname pvn@margolus_size64.bit_exact_reversal\endcsname{10}
\expandafter\def\csname pvtrue@margolus_size64.bit_exact_reversal\endcsname{10}
\expandafter\def\csname pv@margolus_size64.endpoint_recovered\endcsname{10/10}
\expandafter\def\csname pvfrac@margolus_size64.endpoint_recovered\endcsname{1}
\expandafter\def\csname pvn@margolus_size64.endpoint_recovered\endcsname{10}
\expandafter\def\csname pvtrue@margolus_size64.endpoint_recovered\endcsname{10}
\expandafter\def\csname pv@margolus_size64.entropy_conserved\endcsname{10/10}
\expandafter\def\csname pvfrac@margolus_size64.entropy_conserved\endcsname{1}
\expandafter\def\csname pvn@margolus_size64.entropy_conserved\endcsname{10}
\expandafter\def\csname pvtrue@margolus_size64.entropy_conserved\endcsname{10}
\expandafter\def\csname pv@margolus_size64.hamming_weight\endcsname{820.6}
\expandafter\def\csname pvci@margolus_size64.hamming_weight\endcsname{820.6 (35.95)}
\expandafter\def\csname pvsd@margolus_size64.hamming_weight\endcsname{35.95}
\expandafter\def\csname pvn@margolus_size64.hamming_weight\endcsname{10}
\expandafter\def\csname pvmin@margolus_size64.hamming_weight\endcsname{752}
\expandafter\def\csname pvmax@margolus_size64.hamming_weight\endcsname{863}
\expandafter\def\csname pv@margolus_size64.midpoint_match\endcsname{10/10}
\expandafter\def\csname pvfrac@margolus_size64.midpoint_match\endcsname{1}
\expandafter\def\csname pvn@margolus_size64.midpoint_match\endcsname{10}
\expandafter\def\csname pvtrue@margolus_size64.midpoint_match\endcsname{10}
\expandafter\def\csname pv@margolus_steps1.x\endcsname{1}
\expandafter\def\csname pv@margolus_steps1.peak_ws_mb\endcsname{33.17}
\expandafter\def\csname pv@margolus_steps1.peak_ws_sd_mb\endcsname{0.1305}
\expandafter\def\csname pv@margolus_steps1.elapsed_s\endcsname{0.2362}
\expandafter\def\csname pv@margolus_steps1.bit_exact_reversal\endcsname{5/5}
\expandafter\def\csname pvfrac@margolus_steps1.bit_exact_reversal\endcsname{1}
\expandafter\def\csname pvn@margolus_steps1.bit_exact_reversal\endcsname{5}
\expandafter\def\csname pvtrue@margolus_steps1.bit_exact_reversal\endcsname{5}
\expandafter\def\csname pv@margolus_steps1.endpoint_recovered\endcsname{5/5}
\expandafter\def\csname pvfrac@margolus_steps1.endpoint_recovered\endcsname{1}
\expandafter\def\csname pvn@margolus_steps1.endpoint_recovered\endcsname{5}
\expandafter\def\csname pvtrue@margolus_steps1.endpoint_recovered\endcsname{5}
\expandafter\def\csname pv@margolus_steps1.entropy_conserved\endcsname{5/5}
\expandafter\def\csname pvfrac@margolus_steps1.entropy_conserved\endcsname{1}
\expandafter\def\csname pvn@margolus_steps1.entropy_conserved\endcsname{5}
\expandafter\def\csname pvtrue@margolus_steps1.entropy_conserved\endcsname{5}
\expandafter\def\csname pv@margolus_steps1.hamming_weight\endcsname{200.6}
\expandafter\def\csname pvci@margolus_steps1.hamming_weight\endcsname{200.6 (7.925)}
\expandafter\def\csname pvsd@margolus_steps1.hamming_weight\endcsname{7.925}
\expandafter\def\csname pvn@margolus_steps1.hamming_weight\endcsname{5}
\expandafter\def\csname pvmin@margolus_steps1.hamming_weight\endcsname{190}
\expandafter\def\csname pvmax@margolus_steps1.hamming_weight\endcsname{212}
\expandafter\def\csname pv@margolus_steps1.midpoint_match\endcsname{5/5}
\expandafter\def\csname pvfrac@margolus_steps1.midpoint_match\endcsname{1}
\expandafter\def\csname pvn@margolus_steps1.midpoint_match\endcsname{5}
\expandafter\def\csname pvtrue@margolus_steps1.midpoint_match\endcsname{5}
\expandafter\def\csname pv@margolus_steps2.x\endcsname{2}
\expandafter\def\csname pv@margolus_steps2.peak_ws_mb\endcsname{33.25}
\expandafter\def\csname pv@margolus_steps2.peak_ws_sd_mb\endcsname{0.1197}
\expandafter\def\csname pv@margolus_steps2.elapsed_s\endcsname{0.2447}
\expandafter\def\csname pv@margolus_steps2.bit_exact_reversal\endcsname{5/5}
\expandafter\def\csname pvfrac@margolus_steps2.bit_exact_reversal\endcsname{1}
\expandafter\def\csname pvn@margolus_steps2.bit_exact_reversal\endcsname{5}
\expandafter\def\csname pvtrue@margolus_steps2.bit_exact_reversal\endcsname{5}
\expandafter\def\csname pv@margolus_steps2.endpoint_recovered\endcsname{5/5}
\expandafter\def\csname pvfrac@margolus_steps2.endpoint_recovered\endcsname{1}
\expandafter\def\csname pvn@margolus_steps2.endpoint_recovered\endcsname{5}
\expandafter\def\csname pvtrue@margolus_steps2.endpoint_recovered\endcsname{5}
\expandafter\def\csname pv@margolus_steps2.entropy_conserved\endcsname{5/5}
\expandafter\def\csname pvfrac@margolus_steps2.entropy_conserved\endcsname{1}
\expandafter\def\csname pvn@margolus_steps2.entropy_conserved\endcsname{5}
\expandafter\def\csname pvtrue@margolus_steps2.entropy_conserved\endcsname{5}
\expandafter\def\csname pv@margolus_steps2.hamming_weight\endcsname{200.6}
\expandafter\def\csname pvci@margolus_steps2.hamming_weight\endcsname{200.6 (7.925)}
\expandafter\def\csname pvsd@margolus_steps2.hamming_weight\endcsname{7.925}
\expandafter\def\csname pvn@margolus_steps2.hamming_weight\endcsname{5}
\expandafter\def\csname pvmin@margolus_steps2.hamming_weight\endcsname{190}
\expandafter\def\csname pvmax@margolus_steps2.hamming_weight\endcsname{212}
\expandafter\def\csname pv@margolus_steps2.midpoint_match\endcsname{5/5}
\expandafter\def\csname pvfrac@margolus_steps2.midpoint_match\endcsname{1}
\expandafter\def\csname pvn@margolus_steps2.midpoint_match\endcsname{5}
\expandafter\def\csname pvtrue@margolus_steps2.midpoint_match\endcsname{5}
\expandafter\def\csname pv@margolus_steps3.x\endcsname{3}
\expandafter\def\csname pv@margolus_steps3.peak_ws_mb\endcsname{33.34}
\expandafter\def\csname pv@margolus_steps3.peak_ws_sd_mb\endcsname{0.1334}
\expandafter\def\csname pv@margolus_steps3.elapsed_s\endcsname{0.238}
\expandafter\def\csname pv@margolus_steps3.bit_exact_reversal\endcsname{5/5}
\expandafter\def\csname pvfrac@margolus_steps3.bit_exact_reversal\endcsname{1}
\expandafter\def\csname pvn@margolus_steps3.bit_exact_reversal\endcsname{5}
\expandafter\def\csname pvtrue@margolus_steps3.bit_exact_reversal\endcsname{5}
\expandafter\def\csname pv@margolus_steps3.endpoint_recovered\endcsname{5/5}
\expandafter\def\csname pvfrac@margolus_steps3.endpoint_recovered\endcsname{1}
\expandafter\def\csname pvn@margolus_steps3.endpoint_recovered\endcsname{5}
\expandafter\def\csname pvtrue@margolus_steps3.endpoint_recovered\endcsname{5}
\expandafter\def\csname pv@margolus_steps3.entropy_conserved\endcsname{5/5}
\expandafter\def\csname pvfrac@margolus_steps3.entropy_conserved\endcsname{1}
\expandafter\def\csname pvn@margolus_steps3.entropy_conserved\endcsname{5}
\expandafter\def\csname pvtrue@margolus_steps3.entropy_conserved\endcsname{5}
\expandafter\def\csname pv@margolus_steps3.hamming_weight\endcsname{200.6}
\expandafter\def\csname pvci@margolus_steps3.hamming_weight\endcsname{200.6 (7.925)}
\expandafter\def\csname pvsd@margolus_steps3.hamming_weight\endcsname{7.925}
\expandafter\def\csname pvn@margolus_steps3.hamming_weight\endcsname{5}
\expandafter\def\csname pvmin@margolus_steps3.hamming_weight\endcsname{190}
\expandafter\def\csname pvmax@margolus_steps3.hamming_weight\endcsname{212}
\expandafter\def\csname pv@margolus_steps3.midpoint_match\endcsname{5/5}
\expandafter\def\csname pvfrac@margolus_steps3.midpoint_match\endcsname{1}
\expandafter\def\csname pvn@margolus_steps3.midpoint_match\endcsname{5}
\expandafter\def\csname pvtrue@margolus_steps3.midpoint_match\endcsname{5}
\expandafter\def\csname pv@margolus_steps7.x\endcsname{7}
\expandafter\def\csname pv@margolus_steps7.peak_ws_mb\endcsname{33.24}
\expandafter\def\csname pv@margolus_steps7.peak_ws_sd_mb\endcsname{0.1594}
\expandafter\def\csname pv@margolus_steps7.elapsed_s\endcsname{0.2692}
\expandafter\def\csname pv@margolus_steps7.bit_exact_reversal\endcsname{5/5}
\expandafter\def\csname pvfrac@margolus_steps7.bit_exact_reversal\endcsname{1}
\expandafter\def\csname pvn@margolus_steps7.bit_exact_reversal\endcsname{5}
\expandafter\def\csname pvtrue@margolus_steps7.bit_exact_reversal\endcsname{5}
\expandafter\def\csname pv@margolus_steps7.endpoint_recovered\endcsname{5/5}
\expandafter\def\csname pvfrac@margolus_steps7.endpoint_recovered\endcsname{1}
\expandafter\def\csname pvn@margolus_steps7.endpoint_recovered\endcsname{5}
\expandafter\def\csname pvtrue@margolus_steps7.endpoint_recovered\endcsname{5}
\expandafter\def\csname pv@margolus_steps7.entropy_conserved\endcsname{5/5}
\expandafter\def\csname pvfrac@margolus_steps7.entropy_conserved\endcsname{1}
\expandafter\def\csname pvn@margolus_steps7.entropy_conserved\endcsname{5}
\expandafter\def\csname pvtrue@margolus_steps7.entropy_conserved\endcsname{5}
\expandafter\def\csname pv@margolus_steps7.hamming_weight\endcsname{200.6}
\expandafter\def\csname pvci@margolus_steps7.hamming_weight\endcsname{200.6 (7.925)}
\expandafter\def\csname pvsd@margolus_steps7.hamming_weight\endcsname{7.925}
\expandafter\def\csname pvn@margolus_steps7.hamming_weight\endcsname{5}
\expandafter\def\csname pvmin@margolus_steps7.hamming_weight\endcsname{190}
\expandafter\def\csname pvmax@margolus_steps7.hamming_weight\endcsname{212}
\expandafter\def\csname pv@margolus_steps7.midpoint_match\endcsname{5/5}
\expandafter\def\csname pvfrac@margolus_steps7.midpoint_match\endcsname{1}
\expandafter\def\csname pvn@margolus_steps7.midpoint_match\endcsname{5}
\expandafter\def\csname pvtrue@margolus_steps7.midpoint_match\endcsname{5}
\expandafter\def\csname pv@margolus_steps16.x\endcsname{16}
\expandafter\def\csname pv@margolus_steps16.peak_ws_mb\endcsname{33.33}
\expandafter\def\csname pv@margolus_steps16.peak_ws_sd_mb\endcsname{0.1091}
\expandafter\def\csname pv@margolus_steps16.elapsed_s\endcsname{0.2764}
\expandafter\def\csname pv@margolus_steps16.bit_exact_reversal\endcsname{5/5}
\expandafter\def\csname pvfrac@margolus_steps16.bit_exact_reversal\endcsname{1}
\expandafter\def\csname pvn@margolus_steps16.bit_exact_reversal\endcsname{5}
\expandafter\def\csname pvtrue@margolus_steps16.bit_exact_reversal\endcsname{5}
\expandafter\def\csname pv@margolus_steps16.endpoint_recovered\endcsname{5/5}
\expandafter\def\csname pvfrac@margolus_steps16.endpoint_recovered\endcsname{1}
\expandafter\def\csname pvn@margolus_steps16.endpoint_recovered\endcsname{5}
\expandafter\def\csname pvtrue@margolus_steps16.endpoint_recovered\endcsname{5}
\expandafter\def\csname pv@margolus_steps16.entropy_conserved\endcsname{5/5}
\expandafter\def\csname pvfrac@margolus_steps16.entropy_conserved\endcsname{1}
\expandafter\def\csname pvn@margolus_steps16.entropy_conserved\endcsname{5}
\expandafter\def\csname pvtrue@margolus_steps16.entropy_conserved\endcsname{5}
\expandafter\def\csname pv@margolus_steps16.hamming_weight\endcsname{200.6}
\expandafter\def\csname pvci@margolus_steps16.hamming_weight\endcsname{200.6 (7.925)}
\expandafter\def\csname pvsd@margolus_steps16.hamming_weight\endcsname{7.925}
\expandafter\def\csname pvn@margolus_steps16.hamming_weight\endcsname{5}
\expandafter\def\csname pvmin@margolus_steps16.hamming_weight\endcsname{190}
\expandafter\def\csname pvmax@margolus_steps16.hamming_weight\endcsname{212}
\expandafter\def\csname pv@margolus_steps16.midpoint_match\endcsname{5/5}
\expandafter\def\csname pvfrac@margolus_steps16.midpoint_match\endcsname{1}
\expandafter\def\csname pvn@margolus_steps16.midpoint_match\endcsname{5}
\expandafter\def\csname pvtrue@margolus_steps16.midpoint_match\endcsname{5}
\expandafter\def\csname pv@margolus_steps100.x\endcsname{100}
\expandafter\def\csname pv@margolus_steps100.peak_ws_mb\endcsname{33.29}
\expandafter\def\csname pv@margolus_steps100.peak_ws_sd_mb\endcsname{0.1697}
\expandafter\def\csname pv@margolus_steps100.elapsed_s\endcsname{0.3066}
\expandafter\def\csname pv@margolus_steps100.bit_exact_reversal\endcsname{5/5}
\expandafter\def\csname pvfrac@margolus_steps100.bit_exact_reversal\endcsname{1}
\expandafter\def\csname pvn@margolus_steps100.bit_exact_reversal\endcsname{5}
\expandafter\def\csname pvtrue@margolus_steps100.bit_exact_reversal\endcsname{5}
\expandafter\def\csname pv@margolus_steps100.endpoint_recovered\endcsname{5/5}
\expandafter\def\csname pvfrac@margolus_steps100.endpoint_recovered\endcsname{1}
\expandafter\def\csname pvn@margolus_steps100.endpoint_recovered\endcsname{5}
\expandafter\def\csname pvtrue@margolus_steps100.endpoint_recovered\endcsname{5}
\expandafter\def\csname pv@margolus_steps100.entropy_conserved\endcsname{5/5}
\expandafter\def\csname pvfrac@margolus_steps100.entropy_conserved\endcsname{1}
\expandafter\def\csname pvn@margolus_steps100.entropy_conserved\endcsname{5}
\expandafter\def\csname pvtrue@margolus_steps100.entropy_conserved\endcsname{5}
\expandafter\def\csname pv@margolus_steps100.hamming_weight\endcsname{200.6}
\expandafter\def\csname pvci@margolus_steps100.hamming_weight\endcsname{200.6 (7.925)}
\expandafter\def\csname pvsd@margolus_steps100.hamming_weight\endcsname{7.925}
\expandafter\def\csname pvn@margolus_steps100.hamming_weight\endcsname{5}
\expandafter\def\csname pvmin@margolus_steps100.hamming_weight\endcsname{190}
\expandafter\def\csname pvmax@margolus_steps100.hamming_weight\endcsname{212}
\expandafter\def\csname pv@margolus_steps100.midpoint_match\endcsname{5/5}
\expandafter\def\csname pvfrac@margolus_steps100.midpoint_match\endcsname{1}
\expandafter\def\csname pvn@margolus_steps100.midpoint_match\endcsname{5}
\expandafter\def\csname pvtrue@margolus_steps100.midpoint_match\endcsname{5}
\expandafter\def\csname pv@margolus_steps500.x\endcsname{500}
\expandafter\def\csname pv@margolus_steps500.peak_ws_mb\endcsname{33.55}
\expandafter\def\csname pv@margolus_steps500.peak_ws_sd_mb\endcsname{0.1238}
\expandafter\def\csname pv@margolus_steps500.elapsed_s\endcsname{0.6433}
\expandafter\def\csname pv@margolus_steps500.bit_exact_reversal\endcsname{5/5}
\expandafter\def\csname pvfrac@margolus_steps500.bit_exact_reversal\endcsname{1}
\expandafter\def\csname pvn@margolus_steps500.bit_exact_reversal\endcsname{5}
\expandafter\def\csname pvtrue@margolus_steps500.bit_exact_reversal\endcsname{5}
\expandafter\def\csname pv@margolus_steps500.endpoint_recovered\endcsname{5/5}
\expandafter\def\csname pvfrac@margolus_steps500.endpoint_recovered\endcsname{1}
\expandafter\def\csname pvn@margolus_steps500.endpoint_recovered\endcsname{5}
\expandafter\def\csname pvtrue@margolus_steps500.endpoint_recovered\endcsname{5}
\expandafter\def\csname pv@margolus_steps500.entropy_conserved\endcsname{5/5}
\expandafter\def\csname pvfrac@margolus_steps500.entropy_conserved\endcsname{1}
\expandafter\def\csname pvn@margolus_steps500.entropy_conserved\endcsname{5}
\expandafter\def\csname pvtrue@margolus_steps500.entropy_conserved\endcsname{5}
\expandafter\def\csname pv@margolus_steps500.hamming_weight\endcsname{200.6}
\expandafter\def\csname pvci@margolus_steps500.hamming_weight\endcsname{200.6 (7.925)}
\expandafter\def\csname pvsd@margolus_steps500.hamming_weight\endcsname{7.925}
\expandafter\def\csname pvn@margolus_steps500.hamming_weight\endcsname{5}
\expandafter\def\csname pvmin@margolus_steps500.hamming_weight\endcsname{190}
\expandafter\def\csname pvmax@margolus_steps500.hamming_weight\endcsname{212}
\expandafter\def\csname pv@margolus_steps500.midpoint_match\endcsname{5/5}
\expandafter\def\csname pvfrac@margolus_steps500.midpoint_match\endcsname{1}
\expandafter\def\csname pvn@margolus_steps500.midpoint_match\endcsname{5}
\expandafter\def\csname pvtrue@margolus_steps500.midpoint_match\endcsname{5}
\expandafter\def\csname pv@topo_size16.x\endcsname{16}
\expandafter\def\csname pv@topo_size16.peak_ws_mb\endcsname{33.07}
\expandafter\def\csname pv@topo_size16.peak_ws_sd_mb\endcsname{0.1872}
\expandafter\def\csname pv@topo_size16.elapsed_s\endcsname{0.2861}
\expandafter\def\csname pv@topo_size16.accepted_flips\endcsname{76.4}
\expandafter\def\csname pvci@topo_size16.accepted_flips\endcsname{76.4 (7.382)}
\expandafter\def\csname pvsd@topo_size16.accepted_flips\endcsname{7.382}
\expandafter\def\csname pvn@topo_size16.accepted_flips\endcsname{10}
\expandafter\def\csname pvmin@topo_size16.accepted_flips\endcsname{66}
\expandafter\def\csname pvmax@topo_size16.accepted_flips\endcsname{88}
\expandafter\def\csname pv@topo_size16.chi_final\endcsname{0}
\expandafter\def\csname pvci@topo_size16.chi_final\endcsname{0 (0)}
\expandafter\def\csname pvsd@topo_size16.chi_final\endcsname{0}
\expandafter\def\csname pvn@topo_size16.chi_final\endcsname{10}
\expandafter\def\csname pvmin@topo_size16.chi_final\endcsname{0}
\expandafter\def\csname pvmax@topo_size16.chi_final\endcsname{0}
\expandafter\def\csname pv@topo_size16.chi_initial\endcsname{0}
\expandafter\def\csname pvci@topo_size16.chi_initial\endcsname{0 (0)}
\expandafter\def\csname pvsd@topo_size16.chi_initial\endcsname{0}
\expandafter\def\csname pvn@topo_size16.chi_initial\endcsname{10}
\expandafter\def\csname pvmin@topo_size16.chi_initial\endcsname{0}
\expandafter\def\csname pvmax@topo_size16.chi_initial\endcsname{0}
\expandafter\def\csname pv@topo_size16.rejected_flips\endcsname{123.6}
\expandafter\def\csname pvci@topo_size16.rejected_flips\endcsname{123.6 (7.382)}
\expandafter\def\csname pvsd@topo_size16.rejected_flips\endcsname{7.382}
\expandafter\def\csname pvn@topo_size16.rejected_flips\endcsname{10}
\expandafter\def\csname pvmin@topo_size16.rejected_flips\endcsname{112}
\expandafter\def\csname pvmax@topo_size16.rejected_flips\endcsname{134}
\expandafter\def\csname pv@topo_size16.topology_locked\endcsname{10/10}
\expandafter\def\csname pvfrac@topo_size16.topology_locked\endcsname{1}
\expandafter\def\csname pvn@topo_size16.topology_locked\endcsname{10}
\expandafter\def\csname pvtrue@topo_size16.topology_locked\endcsname{10}
\expandafter\def\csname pv@topo_size32.x\endcsname{32}
\expandafter\def\csname pv@topo_size32.peak_ws_mb\endcsname{33.07}
\expandafter\def\csname pv@topo_size32.peak_ws_sd_mb\endcsname{0.1152}
\expandafter\def\csname pv@topo_size32.elapsed_s\endcsname{0.2845}
\expandafter\def\csname pv@topo_size32.accepted_flips\endcsname{35.1}
\expandafter\def\csname pvci@topo_size32.accepted_flips\endcsname{35.1 (5.174)}
\expandafter\def\csname pvsd@topo_size32.accepted_flips\endcsname{5.174}
\expandafter\def\csname pvn@topo_size32.accepted_flips\endcsname{10}
\expandafter\def\csname pvmin@topo_size32.accepted_flips\endcsname{25}
\expandafter\def\csname pvmax@topo_size32.accepted_flips\endcsname{45}
\expandafter\def\csname pv@topo_size32.chi_final\endcsname{0}
\expandafter\def\csname pvci@topo_size32.chi_final\endcsname{0 (0)}
\expandafter\def\csname pvsd@topo_size32.chi_final\endcsname{0}
\expandafter\def\csname pvn@topo_size32.chi_final\endcsname{10}
\expandafter\def\csname pvmin@topo_size32.chi_final\endcsname{0}
\expandafter\def\csname pvmax@topo_size32.chi_final\endcsname{0}
\expandafter\def\csname pv@topo_size32.chi_initial\endcsname{0}
\expandafter\def\csname pvci@topo_size32.chi_initial\endcsname{0 (0)}
\expandafter\def\csname pvsd@topo_size32.chi_initial\endcsname{0}
\expandafter\def\csname pvn@topo_size32.chi_initial\endcsname{10}
\expandafter\def\csname pvmin@topo_size32.chi_initial\endcsname{0}
\expandafter\def\csname pvmax@topo_size32.chi_initial\endcsname{0}
\expandafter\def\csname pv@topo_size32.rejected_flips\endcsname{164.9}
\expandafter\def\csname pvci@topo_size32.rejected_flips\endcsname{164.9 (5.174)}
\expandafter\def\csname pvsd@topo_size32.rejected_flips\endcsname{5.174}
\expandafter\def\csname pvn@topo_size32.rejected_flips\endcsname{10}
\expandafter\def\csname pvmin@topo_size32.rejected_flips\endcsname{155}
\expandafter\def\csname pvmax@topo_size32.rejected_flips\endcsname{175}
\expandafter\def\csname pv@topo_size32.topology_locked\endcsname{10/10}
\expandafter\def\csname pvfrac@topo_size32.topology_locked\endcsname{1}
\expandafter\def\csname pvn@topo_size32.topology_locked\endcsname{10}
\expandafter\def\csname pvtrue@topo_size32.topology_locked\endcsname{10}
\expandafter\def\csname pv@topo_size64.x\endcsname{64}
\expandafter\def\csname pv@topo_size64.peak_ws_mb\endcsname{33.19}
\expandafter\def\csname pv@topo_size64.peak_ws_sd_mb\endcsname{0.1005}
\expandafter\def\csname pv@topo_size64.elapsed_s\endcsname{0.3171}
\expandafter\def\csname pv@topo_size64.accepted_flips\endcsname{7.3}
\expandafter\def\csname pvci@topo_size64.accepted_flips\endcsname{7.3 (2.791)}
\expandafter\def\csname pvsd@topo_size64.accepted_flips\endcsname{2.791}
\expandafter\def\csname pvn@topo_size64.accepted_flips\endcsname{10}
\expandafter\def\csname pvmin@topo_size64.accepted_flips\endcsname{3}
\expandafter\def\csname pvmax@topo_size64.accepted_flips\endcsname{12}
\expandafter\def\csname pv@topo_size64.chi_final\endcsname{0}
\expandafter\def\csname pvci@topo_size64.chi_final\endcsname{0 (0)}
\expandafter\def\csname pvsd@topo_size64.chi_final\endcsname{0}
\expandafter\def\csname pvn@topo_size64.chi_final\endcsname{10}
\expandafter\def\csname pvmin@topo_size64.chi_final\endcsname{0}
\expandafter\def\csname pvmax@topo_size64.chi_final\endcsname{0}
\expandafter\def\csname pv@topo_size64.chi_initial\endcsname{0}
\expandafter\def\csname pvci@topo_size64.chi_initial\endcsname{0 (0)}
\expandafter\def\csname pvsd@topo_size64.chi_initial\endcsname{0}
\expandafter\def\csname pvn@topo_size64.chi_initial\endcsname{10}
\expandafter\def\csname pvmin@topo_size64.chi_initial\endcsname{0}
\expandafter\def\csname pvmax@topo_size64.chi_initial\endcsname{0}
\expandafter\def\csname pv@topo_size64.rejected_flips\endcsname{192.7}
\expandafter\def\csname pvci@topo_size64.rejected_flips\endcsname{192.7 (2.791)}
\expandafter\def\csname pvsd@topo_size64.rejected_flips\endcsname{2.791}
\expandafter\def\csname pvn@topo_size64.rejected_flips\endcsname{10}
\expandafter\def\csname pvmin@topo_size64.rejected_flips\endcsname{188}
\expandafter\def\csname pvmax@topo_size64.rejected_flips\endcsname{197}
\expandafter\def\csname pv@topo_size64.topology_locked\endcsname{10/10}
\expandafter\def\csname pvfrac@topo_size64.topology_locked\endcsname{1}
\expandafter\def\csname pvn@topo_size64.topology_locked\endcsname{10}
\expandafter\def\csname pvtrue@topo_size64.topology_locked\endcsname{10}
\expandafter\def\csname pv@topo_steps2000.x\endcsname{2000}
\expandafter\def\csname pv@topo_steps2000.peak_ws_mb\endcsname{33.2}
\expandafter\def\csname pv@topo_steps2000.peak_ws_sd_mb\endcsname{0.0863}
\expandafter\def\csname pv@topo_steps2000.elapsed_s\endcsname{0.4488}
\expandafter\def\csname pv@topo_steps2000.accepted_flips\endcsname{416}
\expandafter\def\csname pvci@topo_steps2000.accepted_flips\endcsname{416 (36.86)}
\expandafter\def\csname pvsd@topo_steps2000.accepted_flips\endcsname{36.86}
\expandafter\def\csname pvn@topo_steps2000.accepted_flips\endcsname{5}
\expandafter\def\csname pvmin@topo_steps2000.accepted_flips\endcsname{388}
\expandafter\def\csname pvmax@topo_steps2000.accepted_flips\endcsname{479}
\expandafter\def\csname pv@topo_steps2000.chi_final\endcsname{0}
\expandafter\def\csname pvci@topo_steps2000.chi_final\endcsname{0 (0)}
\expandafter\def\csname pvsd@topo_steps2000.chi_final\endcsname{0}
\expandafter\def\csname pvn@topo_steps2000.chi_final\endcsname{5}
\expandafter\def\csname pvmin@topo_steps2000.chi_final\endcsname{0}
\expandafter\def\csname pvmax@topo_steps2000.chi_final\endcsname{0}
\expandafter\def\csname pv@topo_steps2000.chi_initial\endcsname{0}
\expandafter\def\csname pvci@topo_steps2000.chi_initial\endcsname{0 (0)}
\expandafter\def\csname pvsd@topo_steps2000.chi_initial\endcsname{0}
\expandafter\def\csname pvn@topo_steps2000.chi_initial\endcsname{5}
\expandafter\def\csname pvmin@topo_steps2000.chi_initial\endcsname{0}
\expandafter\def\csname pvmax@topo_steps2000.chi_initial\endcsname{0}
\expandafter\def\csname pv@topo_steps2000.rejected_flips\endcsname{1584}
\expandafter\def\csname pvci@topo_steps2000.rejected_flips\endcsname{1584 (36.86)}
\expandafter\def\csname pvsd@topo_steps2000.rejected_flips\endcsname{36.86}
\expandafter\def\csname pvn@topo_steps2000.rejected_flips\endcsname{5}
\expandafter\def\csname pvmin@topo_steps2000.rejected_flips\endcsname{1521}
\expandafter\def\csname pvmax@topo_steps2000.rejected_flips\endcsname{1612}
\expandafter\def\csname pv@topo_steps2000.topology_locked\endcsname{5/5}
\expandafter\def\csname pvfrac@topo_steps2000.topology_locked\endcsname{1}
\expandafter\def\csname pvn@topo_steps2000.topology_locked\endcsname{5}
\expandafter\def\csname pvtrue@topo_steps2000.topology_locked\endcsname{5}
\expandafter\def\csname pv@hpp_size32.x\endcsname{32}
\expandafter\def\csname pv@hpp_size32.peak_ws_mb\endcsname{33.33}
\expandafter\def\csname pv@hpp_size32.peak_ws_sd_mb\endcsname{0.18}
\expandafter\def\csname pv@hpp_size32.elapsed_s\endcsname{0.3828}
\expandafter\def\csname pv@hpp_size32.attraction_detected\endcsname{12/20}
\expandafter\def\csname pvfrac@hpp_size32.attraction_detected\endcsname{0.6}
\expandafter\def\csname pvn@hpp_size32.attraction_detected\endcsname{20}
\expandafter\def\csname pvtrue@hpp_size32.attraction_detected\endcsname{12}
\expandafter\def\csname pv@hpp_size32.final_activity\endcsname{1023}
\expandafter\def\csname pvci@hpp_size32.final_activity\endcsname{1023 (28.1)}
\expandafter\def\csname pvsd@hpp_size32.final_activity\endcsname{28.1}
\expandafter\def\csname pvn@hpp_size32.final_activity\endcsname{20}
\expandafter\def\csname pvmin@hpp_size32.final_activity\endcsname{955}
\expandafter\def\csname pvmax@hpp_size32.final_activity\endcsname{1068}
\expandafter\def\csname pv@hpp_size32.mean_net_force\endcsname{0.0311}
\expandafter\def\csname pvci@hpp_size32.mean_net_force\endcsname{0.0311 (0.1048)}
\expandafter\def\csname pvsd@hpp_size32.mean_net_force\endcsname{0.1048}
\expandafter\def\csname pvn@hpp_size32.mean_net_force\endcsname{20}
\expandafter\def\csname pvmin@hpp_size32.mean_net_force\endcsname{-0.1217}
\expandafter\def\csname pvmax@hpp_size32.mean_net_force\endcsname{0.245}
\expandafter\def\csname pv@hpp_size32.steps\endcsname{300}
\expandafter\def\csname pvci@hpp_size32.steps\endcsname{300 (0)}
\expandafter\def\csname pvsd@hpp_size32.steps\endcsname{0}
\expandafter\def\csname pvn@hpp_size32.steps\endcsname{20}
\expandafter\def\csname pvmin@hpp_size32.steps\endcsname{300}
\expandafter\def\csname pvmax@hpp_size32.steps\endcsname{300}
\expandafter\def\csname pv@hpp_size48.x\endcsname{48}
\expandafter\def\csname pv@hpp_size48.peak_ws_mb\endcsname{33.36}
\expandafter\def\csname pv@hpp_size48.peak_ws_sd_mb\endcsname{0.1426}
\expandafter\def\csname pv@hpp_size48.elapsed_s\endcsname{0.371}
\expandafter\def\csname pv@hpp_size48.attraction_detected\endcsname{10/20}
\expandafter\def\csname pvfrac@hpp_size48.attraction_detected\endcsname{0.5}
\expandafter\def\csname pvn@hpp_size48.attraction_detected\endcsname{20}
\expandafter\def\csname pvtrue@hpp_size48.attraction_detected\endcsname{10}
\expandafter\def\csname pv@hpp_size48.final_activity\endcsname{2301}
\expandafter\def\csname pvci@hpp_size48.final_activity\endcsname{2301 (47.71)}
\expandafter\def\csname pvsd@hpp_size48.final_activity\endcsname{47.71}
\expandafter\def\csname pvn@hpp_size48.final_activity\endcsname{20}
\expandafter\def\csname pvmin@hpp_size48.final_activity\endcsname{2157}
\expandafter\def\csname pvmax@hpp_size48.final_activity\endcsname{2371}
\expandafter\def\csname pv@hpp_size48.mean_net_force\endcsname{-0.006}
\expandafter\def\csname pvci@hpp_size48.mean_net_force\endcsname{-0.006 (0.0461)}
\expandafter\def\csname pvsd@hpp_size48.mean_net_force\endcsname{0.0461}
\expandafter\def\csname pvn@hpp_size48.mean_net_force\endcsname{20}
\expandafter\def\csname pvmin@hpp_size48.mean_net_force\endcsname{-0.105}
\expandafter\def\csname pvmax@hpp_size48.mean_net_force\endcsname{0.0683}
\expandafter\def\csname pv@hpp_size48.steps\endcsname{300}
\expandafter\def\csname pvci@hpp_size48.steps\endcsname{300 (0)}
\expandafter\def\csname pvsd@hpp_size48.steps\endcsname{0}
\expandafter\def\csname pvn@hpp_size48.steps\endcsname{20}
\expandafter\def\csname pvmin@hpp_size48.steps\endcsname{300}
\expandafter\def\csname pvmax@hpp_size48.steps\endcsname{300}
\expandafter\def\csname pv@hpp_size64.x\endcsname{64}
\expandafter\def\csname pv@hpp_size64.peak_ws_mb\endcsname{33.4}
\expandafter\def\csname pv@hpp_size64.peak_ws_sd_mb\endcsname{0.0725}
\expandafter\def\csname pv@hpp_size64.elapsed_s\endcsname{0.3694}
\expandafter\def\csname pv@hpp_size64.attraction_detected\endcsname{12/20}
\expandafter\def\csname pvfrac@hpp_size64.attraction_detected\endcsname{0.6}
\expandafter\def\csname pvn@hpp_size64.attraction_detected\endcsname{20}
\expandafter\def\csname pvtrue@hpp_size64.attraction_detected\endcsname{12}
\expandafter\def\csname pv@hpp_size64.final_activity\endcsname{4098}
\expandafter\def\csname pvci@hpp_size64.final_activity\endcsname{4098 (60.62)}
\expandafter\def\csname pvsd@hpp_size64.final_activity\endcsname{60.62}
\expandafter\def\csname pvn@hpp_size64.final_activity\endcsname{20}
\expandafter\def\csname pvmin@hpp_size64.final_activity\endcsname{3935}
\expandafter\def\csname pvmax@hpp_size64.final_activity\endcsname{4182}
\expandafter\def\csname pv@hpp_size64.mean_net_force\endcsname{0.0066}
\expandafter\def\csname pvci@hpp_size64.mean_net_force\endcsname{0.0066 (0.0549)}
\expandafter\def\csname pvsd@hpp_size64.mean_net_force\endcsname{0.0549}
\expandafter\def\csname pvn@hpp_size64.mean_net_force\endcsname{20}
\expandafter\def\csname pvmin@hpp_size64.mean_net_force\endcsname{-0.095}
\expandafter\def\csname pvmax@hpp_size64.mean_net_force\endcsname{0.0967}
\expandafter\def\csname pv@hpp_size64.steps\endcsname{300}
\expandafter\def\csname pvci@hpp_size64.steps\endcsname{300 (0)}
\expandafter\def\csname pvsd@hpp_size64.steps\endcsname{0}
\expandafter\def\csname pvn@hpp_size64.steps\endcsname{20}
\expandafter\def\csname pvmin@hpp_size64.steps\endcsname{300}
\expandafter\def\csname pvmax@hpp_size64.steps\endcsname{300}
\expandafter\def\csname pv@hpp_size96.x\endcsname{96}
\expandafter\def\csname pv@hpp_size96.peak_ws_mb\endcsname{33.55}
\expandafter\def\csname pv@hpp_size96.peak_ws_sd_mb\endcsname{0.0844}
\expandafter\def\csname pv@hpp_size96.elapsed_s\endcsname{0.457}
\expandafter\def\csname pv@hpp_size96.attraction_detected\endcsname{12/20}
\expandafter\def\csname pvfrac@hpp_size96.attraction_detected\endcsname{0.6}
\expandafter\def\csname pvn@hpp_size96.attraction_detected\endcsname{20}
\expandafter\def\csname pvtrue@hpp_size96.attraction_detected\endcsname{12}
\expandafter\def\csname pv@hpp_size96.final_activity\endcsname{9212}
\expandafter\def\csname pvci@hpp_size96.final_activity\endcsname{9212 (76.15)}
\expandafter\def\csname pvsd@hpp_size96.final_activity\endcsname{76.15}
\expandafter\def\csname pvn@hpp_size96.final_activity\endcsname{20}
\expandafter\def\csname pvmin@hpp_size96.final_activity\endcsname{9056}
\expandafter\def\csname pvmax@hpp_size96.final_activity\endcsname{9330}
\expandafter\def\csname pv@hpp_size96.mean_net_force\endcsname{0.0081}
\expandafter\def\csname pvci@hpp_size96.mean_net_force\endcsname{0.0081 (0.0517)}
\expandafter\def\csname pvsd@hpp_size96.mean_net_force\endcsname{0.0517}
\expandafter\def\csname pvn@hpp_size96.mean_net_force\endcsname{20}
\expandafter\def\csname pvmin@hpp_size96.mean_net_force\endcsname{-0.0917}
\expandafter\def\csname pvmax@hpp_size96.mean_net_force\endcsname{0.1}
\expandafter\def\csname pv@hpp_size96.steps\endcsname{300}
\expandafter\def\csname pvci@hpp_size96.steps\endcsname{300 (0)}
\expandafter\def\csname pvsd@hpp_size96.steps\endcsname{0}
\expandafter\def\csname pvn@hpp_size96.steps\endcsname{20}
\expandafter\def\csname pvmin@hpp_size96.steps\endcsname{300}
\expandafter\def\csname pvmax@hpp_size96.steps\endcsname{300}
\expandafter\def\csname pv@train_depth2.x\endcsname{2}
\expandafter\def\csname pv@train_depth2.peak_ws_mb\endcsname{33.51}
\expandafter\def\csname pv@train_depth2.peak_ws_sd_mb\endcsname{0.1092}
\expandafter\def\csname pv@train_depth2.elapsed_s\endcsname{0.4523}
\expandafter\def\csname pv@train_depth2.epochs\endcsname{120}
\expandafter\def\csname pvci@train_depth2.epochs\endcsname{120 (0)}
\expandafter\def\csname pvsd@train_depth2.epochs\endcsname{0}
\expandafter\def\csname pvn@train_depth2.epochs\endcsname{10}
\expandafter\def\csname pvmin@train_depth2.epochs\endcsname{120}
\expandafter\def\csname pvmax@train_depth2.epochs\endcsname{120}
\expandafter\def\csname pv@train_depth2.loss_decreased\endcsname{10/10}
\expandafter\def\csname pvfrac@train_depth2.loss_decreased\endcsname{1}
\expandafter\def\csname pvn@train_depth2.loss_decreased\endcsname{10}
\expandafter\def\csname pvtrue@train_depth2.loss_decreased\endcsname{10}
\expandafter\def\csname pv@train_depth2.loss_first\endcsname{0.4978}
\expandafter\def\csname pvci@train_depth2.loss_first\endcsname{0.4978 (0.022)}
\expandafter\def\csname pvsd@train_depth2.loss_first\endcsname{0.022}
\expandafter\def\csname pvn@train_depth2.loss_first\endcsname{10}
\expandafter\def\csname pvmin@train_depth2.loss_first\endcsname{0.4705}
\expandafter\def\csname pvmax@train_depth2.loss_first\endcsname{0.5417}
\expandafter\def\csname pv@train_depth2.loss_last\endcsname{0.4223}
\expandafter\def\csname pvci@train_depth2.loss_last\endcsname{0.4223 (0.0179)}
\expandafter\def\csname pvsd@train_depth2.loss_last\endcsname{0.0179}
\expandafter\def\csname pvn@train_depth2.loss_last\endcsname{10}
\expandafter\def\csname pvmin@train_depth2.loss_last\endcsname{0.3947}
\expandafter\def\csname pvmax@train_depth2.loss_last\endcsname{0.4522}
\expandafter\def\csname pv@train_depth2.xor_accuracy_pct\endcsname{58.48}
\expandafter\def\csname pvci@train_depth2.xor_accuracy_pct\endcsname{58.48 (2.613)}
\expandafter\def\csname pvsd@train_depth2.xor_accuracy_pct\endcsname{2.613}
\expandafter\def\csname pvn@train_depth2.xor_accuracy_pct\endcsname{10}
\expandafter\def\csname pvmin@train_depth2.xor_accuracy_pct\endcsname{53.5}
\expandafter\def\csname pvmax@train_depth2.xor_accuracy_pct\endcsname{62}
\expandafter\def\csname pv@train_depth4.x\endcsname{4}
\expandafter\def\csname pv@train_depth4.peak_ws_mb\endcsname{33.57}
\expandafter\def\csname pv@train_depth4.peak_ws_sd_mb\endcsname{0.1737}
\expandafter\def\csname pv@train_depth4.elapsed_s\endcsname{0.3232}
\expandafter\def\csname pv@train_depth4.epochs\endcsname{120}
\expandafter\def\csname pvci@train_depth4.epochs\endcsname{120 (0)}
\expandafter\def\csname pvsd@train_depth4.epochs\endcsname{0}
\expandafter\def\csname pvn@train_depth4.epochs\endcsname{10}
\expandafter\def\csname pvmin@train_depth4.epochs\endcsname{120}
\expandafter\def\csname pvmax@train_depth4.epochs\endcsname{120}
\expandafter\def\csname pv@train_depth4.loss_decreased\endcsname{10/10}
\expandafter\def\csname pvfrac@train_depth4.loss_decreased\endcsname{1}
\expandafter\def\csname pvn@train_depth4.loss_decreased\endcsname{10}
\expandafter\def\csname pvtrue@train_depth4.loss_decreased\endcsname{10}
\expandafter\def\csname pv@train_depth4.loss_first\endcsname{0.5007}
\expandafter\def\csname pvci@train_depth4.loss_first\endcsname{0.5007 (0.0332)}
\expandafter\def\csname pvsd@train_depth4.loss_first\endcsname{0.0332}
\expandafter\def\csname pvn@train_depth4.loss_first\endcsname{10}
\expandafter\def\csname pvmin@train_depth4.loss_first\endcsname{0.4595}
\expandafter\def\csname pvmax@train_depth4.loss_first\endcsname{0.569}
\expandafter\def\csname pv@train_depth4.loss_last\endcsname{0.2628}
\expandafter\def\csname pvci@train_depth4.loss_last\endcsname{0.2628 (0.0132)}
\expandafter\def\csname pvsd@train_depth4.loss_last\endcsname{0.0132}
\expandafter\def\csname pvn@train_depth4.loss_last\endcsname{10}
\expandafter\def\csname pvmin@train_depth4.loss_last\endcsname{0.2452}
\expandafter\def\csname pvmax@train_depth4.loss_last\endcsname{0.2876}
\expandafter\def\csname pv@train_depth4.xor_accuracy_pct\endcsname{76.42}
\expandafter\def\csname pvci@train_depth4.xor_accuracy_pct\endcsname{76.42 (1.922)}
\expandafter\def\csname pvsd@train_depth4.xor_accuracy_pct\endcsname{1.922}
\expandafter\def\csname pvn@train_depth4.xor_accuracy_pct\endcsname{10}
\expandafter\def\csname pvmin@train_depth4.xor_accuracy_pct\endcsname{73.5}
\expandafter\def\csname pvmax@train_depth4.xor_accuracy_pct\endcsname{80}
\expandafter\def\csname pv@train_depth8.x\endcsname{8}
\expandafter\def\csname pv@train_depth8.peak_ws_mb\endcsname{33.59}
\expandafter\def\csname pv@train_depth8.peak_ws_sd_mb\endcsname{0.1513}
\expandafter\def\csname pv@train_depth8.elapsed_s\endcsname{0.3862}
\expandafter\def\csname pv@train_depth8.epochs\endcsname{120}
\expandafter\def\csname pvci@train_depth8.epochs\endcsname{120 (0)}
\expandafter\def\csname pvsd@train_depth8.epochs\endcsname{0}
\expandafter\def\csname pvn@train_depth8.epochs\endcsname{10}
\expandafter\def\csname pvmin@train_depth8.epochs\endcsname{120}
\expandafter\def\csname pvmax@train_depth8.epochs\endcsname{120}
\expandafter\def\csname pv@train_depth8.loss_decreased\endcsname{10/10}
\expandafter\def\csname pvfrac@train_depth8.loss_decreased\endcsname{1}
\expandafter\def\csname pvn@train_depth8.loss_decreased\endcsname{10}
\expandafter\def\csname pvtrue@train_depth8.loss_decreased\endcsname{10}
\expandafter\def\csname pv@train_depth8.loss_first\endcsname{0.4975}
\expandafter\def\csname pvci@train_depth8.loss_first\endcsname{0.4975 (0.0636)}
\expandafter\def\csname pvsd@train_depth8.loss_first\endcsname{0.0636}
\expandafter\def\csname pvn@train_depth8.loss_first\endcsname{10}
\expandafter\def\csname pvmin@train_depth8.loss_first\endcsname{0.3844}
\expandafter\def\csname pvmax@train_depth8.loss_first\endcsname{0.6069}
\expandafter\def\csname pv@train_depth8.loss_last\endcsname{0.1012}
\expandafter\def\csname pvci@train_depth8.loss_last\endcsname{0.1012 (0.0128)}
\expandafter\def\csname pvsd@train_depth8.loss_last\endcsname{0.0128}
\expandafter\def\csname pvn@train_depth8.loss_last\endcsname{10}
\expandafter\def\csname pvmin@train_depth8.loss_last\endcsname{0.0839}
\expandafter\def\csname pvmax@train_depth8.loss_last\endcsname{0.1193}
\expandafter\def\csname pv@train_depth8.xor_accuracy_pct\endcsname{94.47}
\expandafter\def\csname pvci@train_depth8.xor_accuracy_pct\endcsname{94.47 (2.225)}
\expandafter\def\csname pvsd@train_depth8.xor_accuracy_pct\endcsname{2.225}
\expandafter\def\csname pvn@train_depth8.xor_accuracy_pct\endcsname{10}
\expandafter\def\csname pvmin@train_depth8.xor_accuracy_pct\endcsname{90.5}
\expandafter\def\csname pvmax@train_depth8.xor_accuracy_pct\endcsname{97.25}
\expandafter\def\csname pv@train_depth12.x\endcsname{12}
\expandafter\def\csname pv@train_depth12.peak_ws_mb\endcsname{33.58}
\expandafter\def\csname pv@train_depth12.peak_ws_sd_mb\endcsname{0.1897}
\expandafter\def\csname pv@train_depth12.elapsed_s\endcsname{0.4904}
\expandafter\def\csname pv@train_depth12.epochs\endcsname{120}
\expandafter\def\csname pvci@train_depth12.epochs\endcsname{120 (0)}
\expandafter\def\csname pvsd@train_depth12.epochs\endcsname{0}
\expandafter\def\csname pvn@train_depth12.epochs\endcsname{10}
\expandafter\def\csname pvmin@train_depth12.epochs\endcsname{120}
\expandafter\def\csname pvmax@train_depth12.epochs\endcsname{120}
\expandafter\def\csname pv@train_depth12.loss_decreased\endcsname{10/10}
\expandafter\def\csname pvfrac@train_depth12.loss_decreased\endcsname{1}
\expandafter\def\csname pvn@train_depth12.loss_decreased\endcsname{10}
\expandafter\def\csname pvtrue@train_depth12.loss_decreased\endcsname{10}
\expandafter\def\csname pv@train_depth12.loss_first\endcsname{0.4864}
\expandafter\def\csname pvci@train_depth12.loss_first\endcsname{0.4864 (0.0941)}
\expandafter\def\csname pvsd@train_depth12.loss_first\endcsname{0.0941}
\expandafter\def\csname pvn@train_depth12.loss_first\endcsname{10}
\expandafter\def\csname pvmin@train_depth12.loss_first\endcsname{0.3575}
\expandafter\def\csname pvmax@train_depth12.loss_first\endcsname{0.6225}
\expandafter\def\csname pv@train_depth12.loss_last\endcsname{0.0771}
\expandafter\def\csname pvci@train_depth12.loss_last\endcsname{0.0771 (0.0209)}
\expandafter\def\csname pvsd@train_depth12.loss_last\endcsname{0.0209}
\expandafter\def\csname pvn@train_depth12.loss_last\endcsname{10}
\expandafter\def\csname pvmin@train_depth12.loss_last\endcsname{0.0538}
\expandafter\def\csname pvmax@train_depth12.loss_last\endcsname{0.1049}
\expandafter\def\csname pv@train_depth12.xor_accuracy_pct\endcsname{96.42}
\expandafter\def\csname pvci@train_depth12.xor_accuracy_pct\endcsname{96.42 (1.607)}
\expandafter\def\csname pvsd@train_depth12.xor_accuracy_pct\endcsname{1.607}
\expandafter\def\csname pvn@train_depth12.xor_accuracy_pct\endcsname{10}
\expandafter\def\csname pvmin@train_depth12.xor_accuracy_pct\endcsname{94}
\expandafter\def\csname pvmax@train_depth12.xor_accuracy_pct\endcsname{98.5}
\expandafter\def\csname pv@train_depth24.x\endcsname{24}
\expandafter\def\csname pv@train_depth24.peak_ws_mb\endcsname{33.52}
\expandafter\def\csname pv@train_depth24.peak_ws_sd_mb\endcsname{0.1508}
\expandafter\def\csname pv@train_depth24.elapsed_s\endcsname{0.6235}
\expandafter\def\csname pv@train_depth24.epochs\endcsname{120}
\expandafter\def\csname pvci@train_depth24.epochs\endcsname{120 (0)}
\expandafter\def\csname pvsd@train_depth24.epochs\endcsname{0}
\expandafter\def\csname pvn@train_depth24.epochs\endcsname{10}
\expandafter\def\csname pvmin@train_depth24.epochs\endcsname{120}
\expandafter\def\csname pvmax@train_depth24.epochs\endcsname{120}
\expandafter\def\csname pv@train_depth24.loss_decreased\endcsname{10/10}
\expandafter\def\csname pvfrac@train_depth24.loss_decreased\endcsname{1}
\expandafter\def\csname pvn@train_depth24.loss_decreased\endcsname{10}
\expandafter\def\csname pvtrue@train_depth24.loss_decreased\endcsname{10}
\expandafter\def\csname pv@train_depth24.loss_first\endcsname{0.5966}
\expandafter\def\csname pvci@train_depth24.loss_first\endcsname{0.5966 (0.2137)}
\expandafter\def\csname pvsd@train_depth24.loss_first\endcsname{0.2137}
\expandafter\def\csname pvn@train_depth24.loss_first\endcsname{10}
\expandafter\def\csname pvmin@train_depth24.loss_first\endcsname{0.3193}
\expandafter\def\csname pvmax@train_depth24.loss_first\endcsname{0.9315}
\expandafter\def\csname pv@train_depth24.loss_last\endcsname{0.0505}
\expandafter\def\csname pvci@train_depth24.loss_last\endcsname{0.0505 (0.0103)}
\expandafter\def\csname pvsd@train_depth24.loss_last\endcsname{0.0103}
\expandafter\def\csname pvn@train_depth24.loss_last\endcsname{10}
\expandafter\def\csname pvmin@train_depth24.loss_last\endcsname{0.0273}
\expandafter\def\csname pvmax@train_depth24.loss_last\endcsname{0.065}
\expandafter\def\csname pv@train_depth24.xor_accuracy_pct\endcsname{97.65}
\expandafter\def\csname pvci@train_depth24.xor_accuracy_pct\endcsname{97.65 (1.259)}
\expandafter\def\csname pvsd@train_depth24.xor_accuracy_pct\endcsname{1.259}
\expandafter\def\csname pvn@train_depth24.xor_accuracy_pct\endcsname{10}
\expandafter\def\csname pvmin@train_depth24.xor_accuracy_pct\endcsname{96.25}
\expandafter\def\csname pvmax@train_depth24.xor_accuracy_pct\endcsname{99.75}
\expandafter\def\csname pv@train_depth48.x\endcsname{48}
\expandafter\def\csname pv@train_depth48.peak_ws_mb\endcsname{33.49}
\expandafter\def\csname pv@train_depth48.peak_ws_sd_mb\endcsname{0.0689}
\expandafter\def\csname pv@train_depth48.elapsed_s\endcsname{0.9027}
\expandafter\def\csname pv@train_depth48.epochs\endcsname{120}
\expandafter\def\csname pvci@train_depth48.epochs\endcsname{120 (0)}
\expandafter\def\csname pvsd@train_depth48.epochs\endcsname{0}
\expandafter\def\csname pvn@train_depth48.epochs\endcsname{10}
\expandafter\def\csname pvmin@train_depth48.epochs\endcsname{120}
\expandafter\def\csname pvmax@train_depth48.epochs\endcsname{120}
\expandafter\def\csname pv@train_depth48.loss_decreased\endcsname{10/10}
\expandafter\def\csname pvfrac@train_depth48.loss_decreased\endcsname{1}
\expandafter\def\csname pvn@train_depth48.loss_decreased\endcsname{10}
\expandafter\def\csname pvtrue@train_depth48.loss_decreased\endcsname{10}
\expandafter\def\csname pv@train_depth48.loss_first\endcsname{3.467}
\expandafter\def\csname pvci@train_depth48.loss_first\endcsname{3.467 (3.331)}
\expandafter\def\csname pvsd@train_depth48.loss_first\endcsname{3.331}
\expandafter\def\csname pvn@train_depth48.loss_first\endcsname{10}
\expandafter\def\csname pvmin@train_depth48.loss_first\endcsname{0.4919}
\expandafter\def\csname pvmax@train_depth48.loss_first\endcsname{10.25}
\expandafter\def\csname pv@train_depth48.loss_last\endcsname{0.0665}
\expandafter\def\csname pvci@train_depth48.loss_last\endcsname{0.0665 (0.0179)}
\expandafter\def\csname pvsd@train_depth48.loss_last\endcsname{0.0179}
\expandafter\def\csname pvn@train_depth48.loss_last\endcsname{10}
\expandafter\def\csname pvmin@train_depth48.loss_last\endcsname{0.0288}
\expandafter\def\csname pvmax@train_depth48.loss_last\endcsname{0.0954}
\expandafter\def\csname pv@train_depth48.xor_accuracy_pct\endcsname{97}
\expandafter\def\csname pvci@train_depth48.xor_accuracy_pct\endcsname{97 (1.275)}
\expandafter\def\csname pvsd@train_depth48.xor_accuracy_pct\endcsname{1.275}
\expandafter\def\csname pvn@train_depth48.xor_accuracy_pct\endcsname{10}
\expandafter\def\csname pvmin@train_depth48.xor_accuracy_pct\endcsname{95.5}
\expandafter\def\csname pvmax@train_depth48.xor_accuracy_pct\endcsname{99.25}
\expandafter\def\csname pv@train_batch100.x\endcsname{100}
\expandafter\def\csname pv@train_batch100.peak_ws_mb\endcsname{33.33}
\expandafter\def\csname pv@train_batch100.peak_ws_sd_mb\endcsname{0.1678}
\expandafter\def\csname pv@train_batch100.elapsed_s\endcsname{0.3914}
\expandafter\def\csname pv@train_batch100.epochs\endcsname{120}
\expandafter\def\csname pvci@train_batch100.epochs\endcsname{120 (0)}
\expandafter\def\csname pvsd@train_batch100.epochs\endcsname{0}
\expandafter\def\csname pvn@train_batch100.epochs\endcsname{5}
\expandafter\def\csname pvmin@train_batch100.epochs\endcsname{120}
\expandafter\def\csname pvmax@train_batch100.epochs\endcsname{120}
\expandafter\def\csname pv@train_batch100.loss_decreased\endcsname{5/5}
\expandafter\def\csname pvfrac@train_batch100.loss_decreased\endcsname{1}
\expandafter\def\csname pvn@train_batch100.loss_decreased\endcsname{5}
\expandafter\def\csname pvtrue@train_batch100.loss_decreased\endcsname{5}
\expandafter\def\csname pv@train_batch100.loss_first\endcsname{0.4362}
\expandafter\def\csname pvci@train_batch100.loss_first\endcsname{0.4362 (0.0743)}
\expandafter\def\csname pvsd@train_batch100.loss_first\endcsname{0.0743}
\expandafter\def\csname pvn@train_batch100.loss_first\endcsname{5}
\expandafter\def\csname pvmin@train_batch100.loss_first\endcsname{0.3489}
\expandafter\def\csname pvmax@train_batch100.loss_first\endcsname{0.554}
\expandafter\def\csname pv@train_batch100.loss_last\endcsname{0.0548}
\expandafter\def\csname pvci@train_batch100.loss_last\endcsname{0.0548 (0.0236)}
\expandafter\def\csname pvsd@train_batch100.loss_last\endcsname{0.0236}
\expandafter\def\csname pvn@train_batch100.loss_last\endcsname{5}
\expandafter\def\csname pvmin@train_batch100.loss_last\endcsname{0.0317}
\expandafter\def\csname pvmax@train_batch100.loss_last\endcsname{0.0925}
\expandafter\def\csname pv@train_batch100.xor_accuracy_pct\endcsname{96.8}
\expandafter\def\csname pvci@train_batch100.xor_accuracy_pct\endcsname{96.8 (2.387)}
\expandafter\def\csname pvsd@train_batch100.xor_accuracy_pct\endcsname{2.387}
\expandafter\def\csname pvn@train_batch100.xor_accuracy_pct\endcsname{5}
\expandafter\def\csname pvmin@train_batch100.xor_accuracy_pct\endcsname{93}
\expandafter\def\csname pvmax@train_batch100.xor_accuracy_pct\endcsname{99}
\expandafter\def\csname pv@train_batch400.x\endcsname{400}
\expandafter\def\csname pv@train_batch400.peak_ws_mb\endcsname{33.68}
\expandafter\def\csname pv@train_batch400.peak_ws_sd_mb\endcsname{0.1663}
\expandafter\def\csname pv@train_batch400.elapsed_s\endcsname{0.4141}
\expandafter\def\csname pv@train_batch400.epochs\endcsname{120}
\expandafter\def\csname pvci@train_batch400.epochs\endcsname{120 (0)}
\expandafter\def\csname pvsd@train_batch400.epochs\endcsname{0}
\expandafter\def\csname pvn@train_batch400.epochs\endcsname{5}
\expandafter\def\csname pvmin@train_batch400.epochs\endcsname{120}
\expandafter\def\csname pvmax@train_batch400.epochs\endcsname{120}
\expandafter\def\csname pv@train_batch400.loss_decreased\endcsname{5/5}
\expandafter\def\csname pvfrac@train_batch400.loss_decreased\endcsname{1}
\expandafter\def\csname pvn@train_batch400.loss_decreased\endcsname{5}
\expandafter\def\csname pvtrue@train_batch400.loss_decreased\endcsname{5}
\expandafter\def\csname pv@train_batch400.loss_first\endcsname{0.4623}
\expandafter\def\csname pvci@train_batch400.loss_first\endcsname{0.4623 (0.0992)}
\expandafter\def\csname pvsd@train_batch400.loss_first\endcsname{0.0992}
\expandafter\def\csname pvn@train_batch400.loss_first\endcsname{5}
\expandafter\def\csname pvmin@train_batch400.loss_first\endcsname{0.3575}
\expandafter\def\csname pvmax@train_batch400.loss_first\endcsname{0.6225}
\expandafter\def\csname pv@train_batch400.loss_last\endcsname{0.0643}
\expandafter\def\csname pvci@train_batch400.loss_last\endcsname{0.0643 (0.0151)}
\expandafter\def\csname pvsd@train_batch400.loss_last\endcsname{0.0151}
\expandafter\def\csname pvn@train_batch400.loss_last\endcsname{5}
\expandafter\def\csname pvmin@train_batch400.loss_last\endcsname{0.0538}
\expandafter\def\csname pvmax@train_batch400.loss_last\endcsname{0.089}
\expandafter\def\csname pv@train_batch400.xor_accuracy_pct\endcsname{96.5}
\expandafter\def\csname pvci@train_batch400.xor_accuracy_pct\endcsname{96.5 (1.311)}
\expandafter\def\csname pvsd@train_batch400.xor_accuracy_pct\endcsname{1.311}
\expandafter\def\csname pvn@train_batch400.xor_accuracy_pct\endcsname{5}
\expandafter\def\csname pvmin@train_batch400.xor_accuracy_pct\endcsname{94.75}
\expandafter\def\csname pvmax@train_batch400.xor_accuracy_pct\endcsname{98}
\expandafter\def\csname pv@train_batch1600.x\endcsname{1600}
\expandafter\def\csname pv@train_batch1600.peak_ws_mb\endcsname{33.67}
\expandafter\def\csname pv@train_batch1600.peak_ws_sd_mb\endcsname{0.141}
\expandafter\def\csname pv@train_batch1600.elapsed_s\endcsname{0.7384}
\expandafter\def\csname pv@train_batch1600.epochs\endcsname{120}
\expandafter\def\csname pvci@train_batch1600.epochs\endcsname{120 (0)}
\expandafter\def\csname pvsd@train_batch1600.epochs\endcsname{0}
\expandafter\def\csname pvn@train_batch1600.epochs\endcsname{5}
\expandafter\def\csname pvmin@train_batch1600.epochs\endcsname{120}
\expandafter\def\csname pvmax@train_batch1600.epochs\endcsname{120}
\expandafter\def\csname pv@train_batch1600.loss_decreased\endcsname{5/5}
\expandafter\def\csname pvfrac@train_batch1600.loss_decreased\endcsname{1}
\expandafter\def\csname pvn@train_batch1600.loss_decreased\endcsname{5}
\expandafter\def\csname pvtrue@train_batch1600.loss_decreased\endcsname{5}
\expandafter\def\csname pv@train_batch1600.loss_first\endcsname{0.462}
\expandafter\def\csname pvci@train_batch1600.loss_first\endcsname{0.462 (0.0838)}
\expandafter\def\csname pvsd@train_batch1600.loss_first\endcsname{0.0838}
\expandafter\def\csname pvn@train_batch1600.loss_first\endcsname{5}
\expandafter\def\csname pvmin@train_batch1600.loss_first\endcsname{0.3592}
\expandafter\def\csname pvmax@train_batch1600.loss_first\endcsname{0.5815}
\expandafter\def\csname pv@train_batch1600.loss_last\endcsname{0.0629}
\expandafter\def\csname pvci@train_batch1600.loss_last\endcsname{0.0629 (0.0225)}
\expandafter\def\csname pvsd@train_batch1600.loss_last\endcsname{0.0225}
\expandafter\def\csname pvn@train_batch1600.loss_last\endcsname{5}
\expandafter\def\csname pvmin@train_batch1600.loss_last\endcsname{0.0403}
\expandafter\def\csname pvmax@train_batch1600.loss_last\endcsname{0.1003}
\expandafter\def\csname pv@train_batch1600.xor_accuracy_pct\endcsname{97.13}
\expandafter\def\csname pvci@train_batch1600.xor_accuracy_pct\endcsname{97.13 (2.264)}
\expandafter\def\csname pvsd@train_batch1600.xor_accuracy_pct\endcsname{2.264}
\expandafter\def\csname pvn@train_batch1600.xor_accuracy_pct\endcsname{5}
\expandafter\def\csname pvmin@train_batch1600.xor_accuracy_pct\endcsname{93.19}
\expandafter\def\csname pvmax@train_batch1600.xor_accuracy_pct\endcsname{98.75}
\expandafter\def\csname pv@train_batch6400.x\endcsname{6400}
\expandafter\def\csname pv@train_batch6400.peak_ws_mb\endcsname{35.77}
\expandafter\def\csname pv@train_batch6400.peak_ws_sd_mb\endcsname{0.1929}
\expandafter\def\csname pv@train_batch6400.elapsed_s\endcsname{2.072}
\expandafter\def\csname pv@train_batch6400.epochs\endcsname{120}
\expandafter\def\csname pvci@train_batch6400.epochs\endcsname{120 (0)}
\expandafter\def\csname pvsd@train_batch6400.epochs\endcsname{0}
\expandafter\def\csname pvn@train_batch6400.epochs\endcsname{5}
\expandafter\def\csname pvmin@train_batch6400.epochs\endcsname{120}
\expandafter\def\csname pvmax@train_batch6400.epochs\endcsname{120}
\expandafter\def\csname pv@train_batch6400.loss_decreased\endcsname{5/5}
\expandafter\def\csname pvfrac@train_batch6400.loss_decreased\endcsname{1}
\expandafter\def\csname pvn@train_batch6400.loss_decreased\endcsname{5}
\expandafter\def\csname pvtrue@train_batch6400.loss_decreased\endcsname{5}
\expandafter\def\csname pv@train_batch6400.loss_first\endcsname{0.4644}
\expandafter\def\csname pvci@train_batch6400.loss_first\endcsname{0.4644 (0.0856)}
\expandafter\def\csname pvsd@train_batch6400.loss_first\endcsname{0.0856}
\expandafter\def\csname pvn@train_batch6400.loss_first\endcsname{5}
\expandafter\def\csname pvmin@train_batch6400.loss_first\endcsname{0.3548}
\expandafter\def\csname pvmax@train_batch6400.loss_first\endcsname{0.587}
\expandafter\def\csname pv@train_batch6400.loss_last\endcsname{0.0625}
\expandafter\def\csname pvci@train_batch6400.loss_last\endcsname{0.0625 (0.016)}
\expandafter\def\csname pvsd@train_batch6400.loss_last\endcsname{0.016}
\expandafter\def\csname pvn@train_batch6400.loss_last\endcsname{5}
\expandafter\def\csname pvmin@train_batch6400.loss_last\endcsname{0.0529}
\expandafter\def\csname pvmax@train_batch6400.loss_last\endcsname{0.0909}
\expandafter\def\csname pv@train_batch6400.xor_accuracy_pct\endcsname{97.07}
\expandafter\def\csname pvci@train_batch6400.xor_accuracy_pct\endcsname{97.07 (1.695)}
\expandafter\def\csname pvsd@train_batch6400.xor_accuracy_pct\endcsname{1.695}
\expandafter\def\csname pvn@train_batch6400.xor_accuracy_pct\endcsname{5}
\expandafter\def\csname pvmin@train_batch6400.xor_accuracy_pct\endcsname{94.31}
\expandafter\def\csname pvmax@train_batch6400.xor_accuracy_pct\endcsname{98.48}
\expandafter\def\csname pv@hpp_se32.value\endcsname{0.0234}
\expandafter\def\csname pvci@hpp_se32.value\endcsname{0.0234 (0)}
\expandafter\def\csname pvsd@hpp_se32.value\endcsname{0}
\expandafter\def\csname pvn@hpp_se32.value\endcsname{1}
\expandafter\def\csname pvmin@hpp_se32.value\endcsname{0.0234}
\expandafter\def\csname pvmax@hpp_se32.value\endcsname{0.0234}
\expandafter\def\csname pv@hpp_t32.value\endcsname{1.326}
\expandafter\def\csname pvci@hpp_t32.value\endcsname{1.326 (0)}
\expandafter\def\csname pvsd@hpp_t32.value\endcsname{0}
\expandafter\def\csname pvn@hpp_t32.value\endcsname{1}
\expandafter\def\csname pvmin@hpp_t32.value\endcsname{1.326}
\expandafter\def\csname pvmax@hpp_t32.value\endcsname{1.326}
\expandafter\def\csname pv@hpp_se48.value\endcsname{0.0103}
\expandafter\def\csname pvci@hpp_se48.value\endcsname{0.0103 (0)}
\expandafter\def\csname pvsd@hpp_se48.value\endcsname{0}
\expandafter\def\csname pvn@hpp_se48.value\endcsname{1}
\expandafter\def\csname pvmin@hpp_se48.value\endcsname{0.0103}
\expandafter\def\csname pvmax@hpp_se48.value\endcsname{0.0103}
\expandafter\def\csname pv@hpp_t48.value\endcsname{-0.5817}
\expandafter\def\csname pvci@hpp_t48.value\endcsname{-0.5817 (0)}
\expandafter\def\csname pvsd@hpp_t48.value\endcsname{0}
\expandafter\def\csname pvn@hpp_t48.value\endcsname{1}
\expandafter\def\csname pvmin@hpp_t48.value\endcsname{-0.5817}
\expandafter\def\csname pvmax@hpp_t48.value\endcsname{-0.5817}
\expandafter\def\csname pv@hpp_se64.value\endcsname{0.0123}
\expandafter\def\csname pvci@hpp_se64.value\endcsname{0.0123 (0)}
\expandafter\def\csname pvsd@hpp_se64.value\endcsname{0}
\expandafter\def\csname pvn@hpp_se64.value\endcsname{1}
\expandafter\def\csname pvmin@hpp_se64.value\endcsname{0.0123}
\expandafter\def\csname pvmax@hpp_se64.value\endcsname{0.0123}
\expandafter\def\csname pv@hpp_t64.value\endcsname{0.5362}
\expandafter\def\csname pvci@hpp_t64.value\endcsname{0.5362 (0)}
\expandafter\def\csname pvsd@hpp_t64.value\endcsname{0}
\expandafter\def\csname pvn@hpp_t64.value\endcsname{1}
\expandafter\def\csname pvmin@hpp_t64.value\endcsname{0.5362}
\expandafter\def\csname pvmax@hpp_t64.value\endcsname{0.5362}
\expandafter\def\csname pv@hpp_se96.value\endcsname{0.0116}
\expandafter\def\csname pvci@hpp_se96.value\endcsname{0.0116 (0)}
\expandafter\def\csname pvsd@hpp_se96.value\endcsname{0}
\expandafter\def\csname pvn@hpp_se96.value\endcsname{1}
\expandafter\def\csname pvmin@hpp_se96.value\endcsname{0.0116}
\expandafter\def\csname pvmax@hpp_se96.value\endcsname{0.0116}
\expandafter\def\csname pv@hpp_t96.value\endcsname{0.6995}
\expandafter\def\csname pvci@hpp_t96.value\endcsname{0.6995 (0)}
\expandafter\def\csname pvsd@hpp_t96.value\endcsname{0}
\expandafter\def\csname pvn@hpp_t96.value\endcsname{1}
\expandafter\def\csname pvmin@hpp_t96.value\endcsname{0.6995}
\expandafter\def\csname pvmax@hpp_t96.value\endcsname{0.6995}
\expandafter\def\csname pv@hpp_positive_all.value\endcsname{46}
\expandafter\def\csname pvci@hpp_positive_all.value\endcsname{46 (0)}
\expandafter\def\csname pvsd@hpp_positive_all.value\endcsname{0}
\expandafter\def\csname pvn@hpp_positive_all.value\endcsname{1}
\expandafter\def\csname pvmin@hpp_positive_all.value\endcsname{46}
\expandafter\def\csname pvmax@hpp_positive_all.value\endcsname{46}
\expandafter\def\csname pv@hpp_all_n.value\endcsname{80}
\expandafter\def\csname pvci@hpp_all_n.value\endcsname{80 (0)}
\expandafter\def\csname pvsd@hpp_all_n.value\endcsname{0}
\expandafter\def\csname pvn@hpp_all_n.value\endcsname{1}
\expandafter\def\csname pvmin@hpp_all_n.value\endcsname{80}
\expandafter\def\csname pvmax@hpp_all_n.value\endcsname{80}
\expandafter\def\csname pv@topo_accept16.value\endcsname{0.382}
\expandafter\def\csname pvci@topo_accept16.value\endcsname{0.382 (0)}
\expandafter\def\csname pvsd@topo_accept16.value\endcsname{0}
\expandafter\def\csname pvn@topo_accept16.value\endcsname{1}
\expandafter\def\csname pvmin@topo_accept16.value\endcsname{0.382}
\expandafter\def\csname pvmax@topo_accept16.value\endcsname{0.382}
\expandafter\def\csname pv@topo_accept32.value\endcsname{0.1755}
\expandafter\def\csname pvci@topo_accept32.value\endcsname{0.1755 (0)}
\expandafter\def\csname pvsd@topo_accept32.value\endcsname{0}
\expandafter\def\csname pvn@topo_accept32.value\endcsname{1}
\expandafter\def\csname pvmin@topo_accept32.value\endcsname{0.1755}
\expandafter\def\csname pvmax@topo_accept32.value\endcsname{0.1755}
\expandafter\def\csname pv@topo_accept64.value\endcsname{0.0365}
\expandafter\def\csname pvci@topo_accept64.value\endcsname{0.0365 (0)}
\expandafter\def\csname pvsd@topo_accept64.value\endcsname{0}
\expandafter\def\csname pvn@topo_accept64.value\endcsname{1}
\expandafter\def\csname pvmin@topo_accept64.value\endcsname{0.0365}
\expandafter\def\csname pvmax@topo_accept64.value\endcsname{0.0365}
\expandafter\def\csname pv@mem_depth_delta_mb.value\endcsname{-0.025}
\expandafter\def\csname pvci@mem_depth_delta_mb.value\endcsname{-0.025 (0)}
\expandafter\def\csname pvsd@mem_depth_delta_mb.value\endcsname{0}
\expandafter\def\csname pvn@mem_depth_delta_mb.value\endcsname{1}
\expandafter\def\csname pvmin@mem_depth_delta_mb.value\endcsname{-0.025}
\expandafter\def\csname pvmax@mem_depth_delta_mb.value\endcsname{-0.025}
\expandafter\def\csname pv@mem_batch_delta_mb.value\endcsname{2.446}
\expandafter\def\csname pvci@mem_batch_delta_mb.value\endcsname{2.446 (0)}
\expandafter\def\csname pvsd@mem_batch_delta_mb.value\endcsname{0}
\expandafter\def\csname pvn@mem_batch_delta_mb.value\endcsname{1}
\expandafter\def\csname pvmin@mem_batch_delta_mb.value\endcsname{2.446}
\expandafter\def\csname pvmax@mem_batch_delta_mb.value\endcsname{2.446}
\expandafter\def\csname pv@mem_depth_ratio.value\endcsname{0.9993}
\expandafter\def\csname pvci@mem_depth_ratio.value\endcsname{0.9993 (0)}
\expandafter\def\csname pvsd@mem_depth_ratio.value\endcsname{0}
\expandafter\def\csname pvn@mem_depth_ratio.value\endcsname{1}
\expandafter\def\csname pvmin@mem_depth_ratio.value\endcsname{0.9993}
\expandafter\def\csname pvmax@mem_depth_ratio.value\endcsname{0.9993}
\expandafter\def\csname pv@mem_depth_abs_mb.value\endcsname{0.025}
\expandafter\def\csname pvci@mem_depth_abs_mb.value\endcsname{0.025 (0)}
\expandafter\def\csname pvsd@mem_depth_abs_mb.value\endcsname{0}
\expandafter\def\csname pvn@mem_depth_abs_mb.value\endcsname{1}
\expandafter\def\csname pvmin@mem_depth_abs_mb.value\endcsname{0.025}
\expandafter\def\csname pvmax@mem_depth_abs_mb.value\endcsname{0.025}
\makeatother

\newcommand{\val}[1]{\ifcsname pv@#1\endcsname\csname pv@#1\endcsname\else\textbf{??#1??}\fi}
\newcommand{\valci}[1]{\ifcsname pvci@#1\endcsname\csname pvci@#1\endcsname\else\textbf{??#1??}\fi}
\newcommand{\valsd}[1]{\ifcsname pvsd@#1\endcsname\csname pvsd@#1\endcsname\else\textbf{??#1??}\fi}
\newcommand{\valn}[1]{\ifcsname pvn@#1\endcsname\csname pvn@#1\endcsname\else\textbf{??#1??}\fi}
\newcommand{\valfrac}[1]{\ifcsname pvfrac@#1\endcsname\csname pvfrac@#1\endcsname\else\textbf{??#1??}\fi}
\newcommand{\valtrue}[1]{\ifcsname pvtrue@#1\endcsname\csname pvtrue@#1\endcsname\else\textbf{??#1??}\fi}

\title{\bfseries Four reversibility claims, re-measured\\[2pt]
\large An independent audit of a reversible-computing suite, with the numbers it did not report}
\author{Konstantin Parashchuk\\\small Independent researcher, Saint Petersburg, Russia\\
\small \texttt{github.com/kostyk348}}
\date{17 November 2026}

\begin{document}
\maketitle

\begin{abstract}
\noindent
A public repository presents four laboratories of reversible, non-floating-point computing: a
Margolus block cellular automaton that walks backwards in time, a lattice whose Euler
characteristic is held fixed by rejecting moves, a bitwise lattice gas said to show an emergent
attraction, and a symplectic network trained with no stored activations. Each laboratory ships a
command that prints a verdict, and each verdict reproduces exactly on this machine, down to the
last digit of the repository's example. Reproducing an example is a weak test, so this paper
replaces each with a predeclared sweep over seeds, sizes, step counts, batch sizes and depths,
sampling peak working set while every process runs. Two claims hold: the block automaton's
inverse is bit-exact for every lattice size and step count tried, including the odd-step phase
case, and the training peak memory is flat in depth (depth 2 and depth 48 differ by
\val{mem_depth_abs_mb.value}\,MB, inside the resolution of the measurement) while it grows with
batch size (\val{mem_batch_delta_mb.value}\,MB from \val{train_batch100.x} to
\val{train_batch6400.x} samples). One claim
weakens: the reported XOR accuracy is the low end of a distribution, not its centre. One claim
fails: across \val{hpp_all_n.value} lattice-gas runs the net force has the wrong sign about half
the time, and no size shows a mean force more than \val{hpp_t32.value} standard errors from
zero. The paper closes with what it did not test, because the gap between an example that runs
and a claim that survives is the point.
\end{abstract}

\noindent\textbf{Keywords:} reversible computing; block cellular automata; lattice gas; adjoint
training; reproducibility; negative results; measurement methodology

\section{Introduction}

Reversible computing has a 60-year argument behind it. Landauer put a floor under bit erasure
\citep{landauer1961}, Bennett showed computation can be arranged so that no information needs to
be destroyed \citep{bennett1973}, Toffoli turned that into a machine model \citep{toffoli1980},
and Fredkin and Toffoli gave the primitive a conservation law to obey \citep{fredkin1982}. The
argument is old enough that a modern demonstration has a low bar to clear: it only has to run.
That is also the trouble. A demonstration that runs is not a measurement, and the distance
between the two is what a reader cannot check from a README.

This paper takes one such suite, the \texttt{timelock-ca} repository \citep{timelockca}, and
treats its four laboratories as four claims to be tested independently. The repository is a fair
target because it is competent: every claim comes with a command, every command prints
machine-readable output, and every printed number reproduced here on the first try. Nothing that
follows is a complaint about broken code. It is a complaint about the size of the evidence.

Two things are being audited and they should be kept apart.

The first is whether a claim is true as stated. A block cellular automaton whose update is a
permutation can be inverted exactly; that is a property of the rule rather than of the test, and
it should hold for any lattice and any number of steps. The same goes for a training scheme that
stores no activations: memory should be flat in depth. These are checkable, and they check out,
which is worth saying plainly.

The second is whether the reported evidence supports the inference the text draws. Here the
audit finds two different failures. A lattice gas sampled with one seed at one size is not
evidence of an emergent force, and when the sweep is run the force turns out to sit inside its
own noise. A topological invariant conserved by rejecting every move that would change it is
conserved by definition, so the interesting quantity becomes the acceptance rate, which the
repository does not report and which collapses as the lattice grows.

\paragraph{Contributions.}
Section~\ref{sec:method} describes a reproducible measurement harness: a Rust binary, an
experiment plan fixed before running, and a ledger mapping every number in this text to a raw
log. Section~\ref{sec:results} reports the four sweeps. Bit-exact reversal holds over five
lattice sizes and seven step counts, including single ticks and an odd midpoint. The
Euler-characteristic lock conserves its charge, but its acceptance rate falls from
\val{topo_accept16.value} at $16\times16$ to \val{topo_accept64.value} at $64\times64$, so at the
larger size the dynamics is nearly frozen and the conservation is nearly vacuous. The
lattice-gas attraction does not survive replication, with \val{hpp_positive_all.value} positive
runs out of \val{hpp_all_n.value}. Activation memory is flat in depth and linear in batch size,
and the reported XOR accuracy is below the mean of the same configuration over ten seeds.
Section~\ref{sec:not-confirmed} lists what this audit did not test.

\section{Subject and provenance}
\label{sec:subject}

The subject is \texttt{kostyk348/timelock-ca} at commit \texttt{c5c94bb}. It is a Python package
with a four-subcommand command-line interface; each subcommand prints a JSON object with a
verdict field. The four laboratories are:

\begin{enumerate}
\item \texttt{margolus}: a partitioned block cellular automaton on a binary lattice, with the
  16-entry billiard-ball permutation table and alternating block offsets, following
  \citet{margolus1984} and \citet{margolus1987book}. It runs $N$ ticks forward and $N$ backward
  and reports whether the lattice returns bit for bit.
\item \texttt{topo}: two rings stamped on a lattice; random single-bit flips are proposed and
  rejected when they would change the Euler characteristic $\chi = V - E + F$, computed in
  integer arithmetic \citep{kong1989}. The verdict is that $\chi$ is unchanged.
\item \texttt{hpp}: an HPP lattice gas after \citet{hardy1976}, with four boolean channels, a
  head-on collision and two obstacle disks. The verdict is that the mean pressure imbalance is
  positive, read as an attractive force between the disks.
\item \texttt{train}: a leapfrog (symplectic) flow trained by walking the adjoint backwards, so
  that no activation is stored, in the spirit of \citet{chen2018neuralode}. The task is XOR of
  the two input signs, and the verdict is an accuracy above 90\%.
\end{enumerate}

The repository's own numbers reproduce here exactly: seed 0 gives a Hamming weight of 51 at
$16\times16$ over 100 ticks; seed 42 gives a loss of 0.3313 falling to 0.0848 with an accuracy of
94.75\% at depth 12; the topological lock accepts 33 of 200 flips with $\chi$ unchanged; the
lattice gas reports a mean net force of 0.0917 at $64\times64$ \citep{timelockca}. Those four
numbers are the repository's claims and are quoted here as claims. Everything else in this paper
is a measurement made for it.

The environment was a single Windows 11 machine with CPython 3.14.6 and NumPy 2.5.1, one process
at a time, no GPU. That is a limitation and Section~\ref{sec:not-confirmed} says so.

\section{Method}
\label{sec:method}

\paragraph{The plan comes first.} Every run is declared in \texttt{experiments/plan.toml} with its
command line, its seeds, its group and a note saying what the run is for. The plan is the
preregistration of this audit: it was written before the numbers were read, and the prose cannot
cite a run that is not in it. Fixing the plan first is the same argument
\citet{nosek2018} make for preregistration, applied to an audit rather than a trial. Runs are spawned directly, with no shell, so a command line in the
plan is the command line that ran.

\paragraph{Peak memory is measured, not asserted.} An O(1)-memory claim is a memory claim, so the
harness polls the child process's peak working set every 12\,ms through
\texttt{K32GetProcessMemoryInfo} and records it per run. Python's interpreter footprint (about
\val{train_depth2.peak_ws_mb}\,MB here) is common to every run, so differences between depths
are the quantity of interest, and Section~\ref{sec:train} reads the table that way.

\paragraph{Statistics.} Numeric metrics are reported as a mean over seeds with the sample standard
deviation, and where a mean is being compared with zero the ratio of the mean to its standard
error is shown instead of a significance verdict, in the spirit of \citet{wasserstein2016}.
That is a normal-approximation interval; for the two-state counts reported below a resampling
interval would be the better instrument \citep{efron1993bootstrap}, and it is named here as the
obvious extension rather than done.
Boolean verdicts are reported as a count out of $n$; a verdict that is true in every run is
written \val{margolus_size16.bit_exact_reversal}, not as ``passes''.

\paragraph{No typed numbers.} The paper contains no numeric literal that a program did not write.
\texttt{papertool ledger} turns \texttt{data/measurements.json} into TeX macros and one ledger row
per value; \texttt{papertool numbers} fails the build if a lookup has no macro or if a decimal
appears in the prose without an allowlist entry saying why it is not a measurement. The allowlist
ships with the artifact, so the exceptions are auditable too. This matters more than it sounds: a
table typed by hand is where an audit quietly becomes a restatement. The practice follows the
reproducibility rules set out by \citet{sandve2013}.

\paragraph{What replication can and cannot show.} A single repository audited on a single machine
is a case study, not a population estimate \citep{baker2016,osc2015}. The claim here is narrow on
purpose: for these four laboratories, on this machine, with these sweeps, here is which sentences
survive. A different machine would shift the timings and the memory baseline, and
Section~\ref{sec:not-confirmed} marks which conclusions depend on that.

\section{Results}
\label{sec:results}

%%NUMBERS-STRICT-BEGIN
\subsection{Block automaton: the inverse is exact}
\label{sec:margolus}

The Margolus step is a permutation of the 16 block states, applied to a checkerboard of
$2\times2$ blocks with the offset alternating between ticks. A permutation has an inverse, so
forward $N$ ticks followed by backward $N$ ticks must return the starting lattice exactly, and
this is a property of the rule rather than of the lattice size. The sweep confirms the property
and, more usefully, it confirms the phase bookkeeping, which is the part that can be wrong: the
backward pass has to start with the opposite offset.

At the largest size, bit-exact reversal held \val{margolus_size64.bit_exact_reversal} runs, and
Hamming weight was conserved in the same runs, as a permutation must conserve it.
Table~\ref{tab:margolus-size} lists all five sizes.

\begin{table}[t]
\centering
\small
\begin{tabular}{lrrrr}
\toprule
lattice & seeds & bit-exact & Hamming weight & $\sigma$ \\
\midrule
$8\times8$   & \valn{margolus_size8.bit_exact_reversal}  & \val{margolus_size8.bit_exact_reversal}  & \val{margolus_size8.hamming_weight}  & \valsd{margolus_size8.hamming_weight} \\
$16\times16$ & \valn{margolus_size16.bit_exact_reversal} & \val{margolus_size16.bit_exact_reversal} & \val{margolus_size16.hamming_weight} & \valsd{margolus_size16.hamming_weight} \\
$32\times32$ & \valn{margolus_size32.bit_exact_reversal} & \val{margolus_size32.bit_exact_reversal} & \val{margolus_size32.hamming_weight} & \valsd{margolus_size32.hamming_weight} \\
$48\times48$ & \valn{margolus_size48.bit_exact_reversal} & \val{margolus_size48.bit_exact_reversal} & \val{margolus_size48.hamming_weight} & \valsd{margolus_size48.hamming_weight} \\
$64\times64$ & \valn{margolus_size64.bit_exact_reversal} & \val{margolus_size64.bit_exact_reversal} & \val{margolus_size64.hamming_weight} & \valsd{margolus_size64.hamming_weight} \\
\bottomrule
\end{tabular}
\caption{Bit-exact reversal by lattice size, 100 ticks forward and back, ten seeds per size.
``Bit-exact'' counts runs in which the final lattice equalled the initial one, bit for bit.
Hamming weight is the initial particle count; the spread is across seeds, not across ticks.}
\label{tab:margolus-size}
\end{table}

The step-count sweep is the one that could have failed and did not. An odd number of ticks leaves
the lattice on the opposite block offset, so the backward pass must invert the phases in the
reverse order; getting that wrong produces a lattice that looks plausible and is wrong. Single
ticks, three ticks, seven ticks with an odd midpoint, and \val{margolus_steps500.x} ticks all returned to the start,
\val{margolus_steps500.bit_exact_reversal} of the runs at the longest step count
(Table~\ref{tab:margolus-steps}). Integer arithmetic is why this is unconditional: no rounding
enters, so no drift accumulates with $N$.

\begin{table}[t]
\centering
\small
\begin{tabular}{lrrr}
\toprule
ticks forward & seeds & bit-exact & Hamming weight \\
\midrule
\val{margolus_steps1.x}   & \valn{margolus_steps1.bit_exact_reversal}   & \val{margolus_steps1.bit_exact_reversal}   & \val{margolus_steps1.hamming_weight} \\
\val{margolus_steps2.x}   & \valn{margolus_steps2.bit_exact_reversal}   & \val{margolus_steps2.bit_exact_reversal}   & \val{margolus_steps2.hamming_weight} \\
\val{margolus_steps3.x}   & \valn{margolus_steps3.bit_exact_reversal}   & \val{margolus_steps3.bit_exact_reversal}   & \val{margolus_steps3.hamming_weight} \\
\val{margolus_steps7.x}   & \valn{margolus_steps7.bit_exact_reversal}   & \val{margolus_steps7.bit_exact_reversal}   & \val{margolus_steps7.hamming_weight} \\
\val{margolus_steps16.x}  & \valn{margolus_steps16.bit_exact_reversal}  & \val{margolus_steps16.bit_exact_reversal}  & \val{margolus_steps16.hamming_weight} \\
\val{margolus_steps100.x} & \valn{margolus_steps100.bit_exact_reversal} & \val{margolus_steps100.bit_exact_reversal} & \val{margolus_steps100.hamming_weight} \\
\val{margolus_steps500.x} & \valn{margolus_steps500.bit_exact_reversal} & \val{margolus_steps500.bit_exact_reversal} & \val{margolus_steps500.hamming_weight} \\
\bottomrule
\end{tabular}
\caption{Reversal as a function of the number of ticks, $32\times32$ lattice, five seeds per row.
The midpoint, where the forward run stops and the backward run starts, is half the tick count
rounded down, so the odd tick counts exercise the case where the two block phases are not in
their starting order.}
\label{tab:margolus-steps}
\end{table}

\subsection{Topological lock: conserved, and increasingly frozen}
\label{sec:topo}

The second laboratory stamps two rings on a lattice and proposes single-bit flips, rejecting any
flip that would change the Euler characteristic. The published verdict is that $\chi$ is
conserved, and it is: $\chi$ ended at \val{topo_size32.chi_final} in every run at every size, and
the initial value was \val{topo_size32.chi_initial}.

The problem is that this cannot fail. Rejecting every move that changes a quantity conserves it
by construction; the mechanism is the rejection sampler of \citet{metropolis1953}, and what a
sampler does is spend accepted moves. Reporting the conserved value without the acceptance rate
reports the definition and hides the dynamics. When the rate is measured, it falls steeply with
lattice area: \val{topo_accept16.value} of proposals accepted at $16\times16$,
\val{topo_accept32.value} at $32\times32$, and \val{topo_accept64.value} at $64\times64$
(Table~\ref{tab:topo}). At the largest size the lattice is barely moving while a conserved charge
certifies that nothing has changed.

\begin{table}[t]
\centering
\small
\begin{tabular}{lrrrr}
\toprule
lattice & accepted & rejected & acceptance rate & locked \\
\midrule
$16\times16$ & \val{topo_size16.accepted_flips} & \val{topo_size16.rejected_flips} & \val{topo_accept16.value} & \val{topo_size16.topology_locked} \\
$32\times32$ & \val{topo_size32.accepted_flips} & \val{topo_size32.rejected_flips} & \val{topo_accept32.value} & \val{topo_size32.topology_locked} \\
$64\times64$ & \val{topo_size64.accepted_flips} & \val{topo_size64.rejected_flips} & \val{topo_accept64.value} & \val{topo_size64.topology_locked} \\
\midrule
$32\times32$, \val{topo_steps2000.x} proposals & \val{topo_steps2000.accepted_flips} & \val{topo_steps2000.rejected_flips} & held & \val{topo_steps2000.topology_locked} \\
\bottomrule
\end{tabular}
\caption{The Euler-characteristic lock at three sizes with 200 proposals each, ten seeds, plus a
ten-times-longer run at the middle size. Accepted and rejected are means over seeds; ``locked''
counts runs whose final $\chi$ equalled the initial one.}
\label{tab:topo}
\end{table}

The longer run separates two explanations. If the lock froze completely the sampler would stop,
and the configuration would never change; instead \val{topo_steps2000.accepted_flips} flips are
accepted out of \val{topo_steps2000.x} proposals, at the same rate as the short run. The lattice
keeps moving at a rate the flip proposal distribution dictates and the invariant then filters.
Whether that motion explores anything is a separate question this audit did not answer.

Nothing here contradicts the repository's sentence, which says the value is locked. What the
audit adds is that the sentence is weaker than it sounds. A checksum recomputed correctly after
every accepted move is not a memory in the sense topological quantum memories use the term, where
the invariant protects a stored state against a local error model \citep{kitaev2003,dennis2002}.
In digital topology the Euler characteristic fixes some global properties of a binary picture and
leaves the shape free \citep{kong1989}: the rings can wander, reach the border and reform while
$\chi$ stays at \val{topo_size64.chi_final}. The repository's README says as much in a
parenthesis, and the measurement agrees with the parenthesis rather than with the headline.

\subsection{Lattice gas: the attraction that was not there}
\label{sec:hpp}

The third laboratory makes a physical claim, and it is the one that did not survive. The HPP
lattice gas \citep{hardy1976} runs for 300 steps with two obstacle disks; the mean pressure
imbalance on the disks is reported as a net force, and a positive mean is read as attraction
between the disks.

At the repository's configuration the single reported seed does give a positive number: the
seed-42 run at $64\times64$ returns a mean net force of 0.0917 and prints a detected attraction
\citep{timelockca}. Repeating that configuration with twenty seeds instead of one gives a mean of
\val{hpp_size64.mean_net_force} with a standard deviation of \valsd{hpp_size64.mean_net_force},
so the standard error is \val{hpp_se64.value} and the mean sits \val{hpp_t64.value} standard
errors from zero. The sign is decided by the seed: \val{hpp_size64.attraction_detected} runs
report a positive force at $64\times64$, and pooled over four sizes the count is
\val{hpp_positive_all.value} out of \val{hpp_all_n.value} (Table~\ref{tab:hpp},
Figure~\ref{fig:hpp}).

\begin{table}[t]
\centering
\small
\begin{tabular}{lrrrrr}
\toprule
lattice & seeds & mean net force & SD & SE & mean/SE \\
\midrule
$32\times32$ & \valn{hpp_size32.mean_net_force} & \val{hpp_size32.mean_net_force} & \valsd{hpp_size32.mean_net_force} & \val{hpp_se32.value} & \val{hpp_t32.value} \\
$48\times48$ & \valn{hpp_size48.mean_net_force} & \val{hpp_size48.mean_net_force} & \valsd{hpp_size48.mean_net_force} & \val{hpp_se48.value} & \val{hpp_t48.value} \\
$64\times64$ & \valn{hpp_size64.mean_net_force} & \val{hpp_size64.mean_net_force} & \valsd{hpp_size64.mean_net_force} & \val{hpp_se64.value} & \val{hpp_t64.value} \\
$96\times96$ & \valn{hpp_size96.mean_net_force} & \val{hpp_size96.mean_net_force} & \valsd{hpp_size96.mean_net_force} & \val{hpp_se96.value} & \val{hpp_t96.value} \\
\bottomrule
\end{tabular}
\caption{Mean net force between the two obstacle disks, 300 steps, twenty seeds per size. SE is
the standard error of the mean; mean/SE expresses the mean in its own standard errors. The
fraction of runs reporting a detected attraction is in the text, since a sign count is not a
force.}
\label{tab:hpp}
\end{table}

\begin{figure}[t]
\centering
\begin{tikzpicture}
\begin{axis}[
  width=0.78\textwidth, height=5.4cm,
  xlabel={lattice edge (cells)}, ylabel={mean net force (lattice units)},
  xmin=24, xmax=104, ymin=-0.10, ymax=0.12,
  grid=major, grid style={gray!25},
  legend pos=north east, legend style={font=\footnotesize},
]
\addplot+[mark=*, thick, error bars/.cd, y dir=both, y explicit]
  table[x=x, y=mean_net_force_mean, y error=mean_net_force_mean_sd, col sep=comma]
  {../data/figures/hpp_size.csv};
\addlegendentry{mean $\pm$ 1 SD over twenty seeds}
\addplot[dashed, black!60] coordinates {(24,0) (104,0)};
\addlegendentry{zero}
\end{axis}
\end{tikzpicture}
\caption{Net force between the obstacle disks against lattice size. Every interval crosses zero
and the four means are small compared with their own spread. The repository's single seed sits
inside this cloud, in either direction; that is what a sample of one buys.}
\label{fig:hpp}
\end{figure}

A force that changes sign with the seed and does not grow with system size is what one expects
from the residual fluctuation of a finite lattice gas rather than from a hydrodynamic
interaction, and picking the seed that produced it is the mechanism \citet{ioannidis2005}
describes for findings that do not replicate. Reproducing the Navier--Stokes limit from a lattice gas requires a specific symmetry
set and collision rule \citep{frisch1986}, and the statistical care with which a lattice
automaton's macroscopic claim has to be established is exactly the subject of \citet{wolfram1986}.
The audited code is a straightforward HPP implementation. The inference from its one printed
number to gravity from pure bits is not supported by \val{hpp_all_n.value} runs, and that is the
finding.

\subsection{Adjoint training: memory flat in depth, linear in batch}
\label{sec:train}

The fourth laboratory makes two claims that are easy to conflate: that the network trains without
storing activations, and that it learns XOR. The first is about memory and is measured directly.
Peak working set was \val{train_depth2.peak_ws_mb}\,MB for a depth-2 flow and
\val{train_depth48.peak_ws_mb}\,MB for a depth-48 flow. The two differ by
\val{mem_depth_abs_mb.value}\,MB in absolute value, and the sign of that difference
(\val{mem_depth_delta_mb.value}\,MB) is the sign of noise: 24 times the depth moves the peak by
less than one part in a thousand (ratio \val{mem_depth_ratio.value}). That is consistent with no per-layer storage at all: stored activations for those shapes
would be megabytes, and checkpointing trades a square-root memory term for recomputation
\citep{griewank2008}, whereas here the depth term is absent.

The batch term is present, as it must be. Going from \val{train_batch100.x} to \val{train_batch6400.x} samples raised the peak by
\val{mem_batch_delta_mb.value}\,MB (Table~\ref{tab:train}), the linear term in the number of
samples that every method pays. The claim that survives is therefore narrower and more useful
than O(1) memory: the method is linear in batch and flat in depth, which is what one wants when
depth is the large number.

\begin{table}[t]
\centering
\small
\begin{tabular}{lrrrrr}
\toprule
depth & seeds & loss (first) & loss (last) & XOR accuracy (\%) & peak WS (MB) \\
\midrule
2  & \valn{train_depth2.xor_accuracy_pct}  & \val{train_depth2.loss_first}  & \val{train_depth2.loss_last}  & \val{train_depth2.xor_accuracy_pct}  & \val{train_depth2.peak_ws_mb} \\
4  & \valn{train_depth4.xor_accuracy_pct}  & \val{train_depth4.loss_first}  & \val{train_depth4.loss_last}  & \val{train_depth4.xor_accuracy_pct}  & \val{train_depth4.peak_ws_mb} \\
8  & \valn{train_depth8.xor_accuracy_pct}  & \val{train_depth8.loss_first}  & \val{train_depth8.loss_last}  & \val{train_depth8.xor_accuracy_pct}  & \val{train_depth8.peak_ws_mb} \\
12 & \valn{train_depth12.xor_accuracy_pct} & \val{train_depth12.loss_first} & \val{train_depth12.loss_last} & \val{train_depth12.xor_accuracy_pct} & \val{train_depth12.peak_ws_mb} \\
24 & \valn{train_depth24.xor_accuracy_pct} & \val{train_depth24.loss_first} & \val{train_depth24.loss_last} & \val{train_depth24.xor_accuracy_pct} & \val{train_depth24.peak_ws_mb} \\
48 & \valn{train_depth48.xor_accuracy_pct} & \val{train_depth48.loss_first} & \val{train_depth48.loss_last} & \val{train_depth48.xor_accuracy_pct} & \val{train_depth48.peak_ws_mb} \\
\bottomrule
\end{tabular}
\caption{Adjoint training against depth, 400 samples, 120 epochs, ten seeds per row. Accuracy is
measured on the training set, because the repository defines no held-out split. Peak working set
is the sampled maximum over the run and includes the Python interpreter (about
\val{train_depth2.peak_ws_mb}\,MB of it), which is why the column is read as differences.}
\label{tab:train}
\end{table}

The accuracy claim behaves differently. The reported 94.75\% at depth 12 with seed 42 is one draw
from a distribution whose mean here is \val{train_depth12.xor_accuracy_pct} with a standard
deviation of \valsd{train_depth12.xor_accuracy_pct}: the reported value sits below the mean and
well inside one standard deviation. The shape of the curve matters more than the level
(Figure~\ref{fig:depth}). Depth 2 settles at \val{train_depth2.xor_accuracy_pct}\% and depth 4 at
\val{train_depth4.xor_accuracy_pct}\%, so the task is not learnable without enough leapfrog
steps, while everything from depth 8 upward lands between
\val{train_depth8.xor_accuracy_pct}\% and \val{train_depth48.xor_accuracy_pct}\%. A capacity story
in which more depth always helps is not what the sweep shows; a threshold near depth 8 is. The
audited command agrees, in its own way: it exits non-zero for every depth-2 and depth-4 run,
because its pass condition is a 90\% accuracy floor, and the harness records the exit code of every
process alongside its output, so that disagreement is visible in the artifact and not only in this
text.

\begin{figure}[t]
\centering
\begin{tikzpicture}
\begin{axis}[
  width=0.62\textwidth, height=5.4cm,
  xlabel={leapfrog depth}, ylabel={XOR accuracy (\%)},
  xmin=0, xmax=52, ymin=50, ymax=100,
  grid=major, grid style={gray!25},
]
\addplot+[mark=*, thick] table[x=x, y=xor_accuracy_pct_mean, col sep=comma]
  {../data/figures/train_depth.csv};
\addplot[dashed, black!60] coordinates {(0,90) (52,90)};
\end{axis}
\end{tikzpicture}
\caption{Training accuracy against depth, mean over ten seeds. The jump happens between depth 4
and depth 8; beyond that the curve is flat within its own spread. The dashed line is the 90\%
threshold the audited command itself uses as its pass condition.}
\label{fig:depth}
\end{figure}

One caveat about the task. The loss pushes orbits toward a target radius of 0.4 or 1.6 depending
on the label, and the readout thresholds the final radius at 1.0. The labels are the XOR of the
two input signs, which is not linearly separable, so the flow has to bend the input square rather
than draw a line through it; that much is real nonlinear work. But the objective and the readout
are two halves of one construction, and a model can score well by learning the radius map without
representing parity at all. The audit measured accuracy, not parity representation, and
Section~\ref{sec:not-confirmed} says so.
%%NUMBERS-STRICT-END

\section{What this paper does not establish}
\label{sec:not-confirmed}

\paragraph{Single machine, one interpreter.} All timings and all memory figures come from one
Windows 11 machine with CPython 3.14.6 and NumPy 2.5.1. The memory conclusions are about
differences between Python processes, in which the interpreter is a large constant term. They do
not establish what a C or Rust implementation of the same scheme would use, and they say nothing
about a GPU.

\paragraph{One repository, one commit.} The subject is a single commit, read as a case study. The
audit says nothing about reversible computing in general, and nothing about other implementations
of the same ideas.

\paragraph{The lattice gas was not characterised exhaustively.} Three hundred steps at four sizes
with twenty seeds each is enough to show that a single-seed positive result is not evidence of a
force. It is not enough to say the force is exactly zero, and it does not identify the mechanism
of the fluctuation, which would need time-averaged pressure profiles and much longer runs.

\paragraph{The topological lock was not tested for ergodicity.} Acceptance rates are reported.
Whether the accepted moves mix, and whether the invariant lets any configuration reach any other
with the same $\chi$, were not measured. The claim that the invariant is a checksum rather than a
memory rests on the acceptance-rate collapse and on the freedom the Euler characteristic leaves in
a binary picture, not on a mixing proof.

\paragraph{The XOR task is not a parity test.} Training accuracy with a radius-based objective is
what the repository computes, and what was measured. Whether the trained flow represents parity
was not tested: that would need a held-out set, a probe for parity structure, and a comparison
against a model allowed to store activations.

\paragraph{No energy measurement.} The motivation for reversible computing is dissipation, and
nothing in this paper measures power. Peak working set is not energy. The repository does not
claim otherwise, and neither does the audit.

\section{Reproducibility statement}

The paper is a package, not a document. \texttt{experiments/plan.toml} declares every run;
\texttt{subjects.lock} pins the audited commit; \texttt{data/measurements.json} holds every
aggregate with the raw value list; \texttt{evidence/raw/} holds one stdout and stderr file per
process; \texttt{data/numbers.tsv} maps every number in this text to the JSON path it came from,
and \texttt{tex/numbers.tex} is generated from that ledger. \texttt{verify\_all.sh} re-runs the
citation checks, the number checks, the ledger freshness check and the PDF build, and refuses to
pass if any of them drifts. The measurement harness is one Rust binary (\texttt{tools/papertool},
standard library only). The audited code is Python, invoked as a black box through its own
command line, so the harness cannot share a bug with the subject.

The one thing the package cannot reproduce on its own is the numbers: they are measurements of a
machine. Re-running the plan on different hardware will shift timings and the memory baseline, so
the paper should be re-read with that in mind. What should survive is the shape of each result,
and the reversal result should survive exactly, because it is integer arithmetic.

\section{Conclusion}

Four claims, three verdicts. Bit-exact reversal is real and unconditional, and the phase
bookkeeping is correct at odd step counts, the case most likely to be wrong. Activation memory is
flat in depth, and the paper's own phrasing should be flat in depth and linear in batch. The
Euler characteristic is conserved by a sampler whose acceptance rate falls to
\val{topo_accept64.value} as the lattice grows, so the conserved quantity is a checksum and the
interesting number is the one the repository does not print. The lattice-gas attraction is not
there: \val{hpp_positive_all.value} positive runs out of \val{hpp_all_n.value}, means inside their
own standard errors, and no size dependence.

That last one is the reason to write this down. A single seed selected a number that reads as
physics; twenty seeds turn it back into a fluctuation. Nothing in the code needed changing for
that to happen, which suggests the fix belongs upstream of the code, in what a demonstration is
allowed to claim before somebody sweeps it.

\section*{Acknowledgements and disclosure}

The author thanks the maintainers of \texttt{timelock-ca} for publishing code whose claims are
checkable; an audit of a repository that hides its numbers would have been impossible. AI
assistance: the measurement harness, the experiment plan, the number ledger, the reference
verification and the drafting were produced in a session with a Hermes Agent instance acting as a
research assistant. Every number in this paper is generated by the artifact from raw logs, and
every reference was resolved by a machine query before it was cited. The checks that enforce both
are in \texttt{verify\_all.sh} and were run by the human author on the published version.

\bibliographystyle{plainnat}
\bibliography{references}

\appendix
\section{Commands}

\begin{verbatim}
bash scripts/fetch_subjects.sh        # clone the subject at the pinned commit
papertool measure --root .            # run the plan, keep logs, aggregate
papertool measure --root . --derive-only
papertool ledger --root .             # measurements -> macros, ledger, figures
papertool bib --root .                # curated sources -> references.bib
papertool cites --root . --check
papertool numbers --root .            # lookups resolve; no unlisted numerals
bash build.sh                         # all of the above, then the PDF
bash verify_all.sh                    # everything, plus the freshness checks\end{verbatim}

\end{document}
